\documentclass[12pt]{report}
\usepackage[portuguese]{babel}
\usepackage[utf8]{inputenc}
%\usepackage[dvips]{graphicx}
\usepackage{graphicx}
\usepackage{enumitem}
\usepackage[normalem]{ulem}
\usepackage{latexsym}
\usepackage{amsmath}
\usepackage{float}
\usepackage{units} % Para frações bonitas
\usepackage{verbatim} % Para que o bloco {comment} funcione
\usepackage[pdfborder={0 0 0}]{hyperref}
\usepackage{wrapfig} % Para o qrcode do editorial
\usepackage[labelformat=empty]{caption} % Para tirar o "Figure:" da caption do QRCode

%\usepackage{showframe} %Para mostrar o frame do arquivo

\newcounter{qcounter}

\setlength{\parindent}{0pt}
\setlength{\parskip}{3ex plus 0.5ex minus 0.5ex}

\addtolength{\voffset}{-2.8cm}
\addtolength{\textheight}{5.5cm}
\addtolength{\hoffset}{-2.3cm}
\addtolength{\textwidth}{4.5cm}
\setlength{\topmargin}{0cm}

%FIXME GAMBIARRA
%\setlength{\evensidemargin}{55pt}

\pagestyle{plain}

% =============== Seção de definições de macros ========================

% Delimita uma seção:
\newenvironment{secao}[1] {
    \framebox[\textwidth] {
        \rule[-1.2ex]{5ex}{5.5ex}
        {\Large\sf #1}
        \hspace{\stretch{1}}
    }
    \phantomsection %Faz o link do pdf funcionar direito
	\addcontentsline{toc}{chapter}{#1}
    \nopagebreak[4]
}{}

%FIXME Isso pode ser resolvido de um jeito melhor, fazendo uma seção (editorial e mandamentos)
% não entrarem no índice usando uma tag, e não outro tipo de environment.
% Veja a primeira linha de editorial.tex para entender melhor.
% Delimita o editorial (que não entra no Índice):
\newenvironment{editorial}[1] {
    \framebox[\textwidth] {
        \rule[-1.2ex]{5ex}{5.5ex}
        {\Large\sf #1}
        \hspace{\stretch{1}}
    }
    \phantomsection %Faz o link do pdf funcionar direito
    %\addcontentsline{toc}{chapter}{#1}
    \nopagebreak[4]
}{
\thispagestyle{empty}
\pagebreak
}


% Delimita uma subseção
\newenvironment{subsecao}[1] {
    \rule[0ex]{2.5ex}{2.5ex}
    {\Large\sf #1}
    \phantomsection %Faz o link do pdf funcionar direito
	\addcontentsline{toc}{section}{#1}
    \nopagebreak[4]
}{}

% Delimita uma subsubseção
\newenvironment{subsubsecao}[1] {
    \rule[0ex]{1ex}{1.5ex}
    {\large\sf #1}
    \phantomsection %Faz o link do pdf funcionar direito
% Próxima linha é responsável por subsubseções não aparecerem no índice!
	%\addcontentsline{toc}{subsection}{#1}
    \nopagebreak[4]
}{}

% Coloca uma figura grande (sem ser quadrinhos)
\newcommand{\figuragrande}[1] {
    \begin{figure}[!htbp]
      \begin{center}
        \includegraphics[width=\textwidth]{img/#1.pdf}
      \end{center}
    \end{figure}
}

% Coloca uma figura menor (sem ser quadrinhos)
\newcommand{\figurapequenainline}[1] {
  \begin{wrapfigure}{r}{0.25\textwidth}
    \vspace{-25pt}
    \begin{center}
      \includegraphics[width=0.25\textwidth]{img/#1.pdf}
    \end{center}
    \vspace{-25pt}
  \end{wrapfigure}
}

% Coloca uma figura menor (sem ser quadrinhos) apertada verticalmente
\newcommand{\figurapequenainlineflexivel}[2] {
  \begin{wrapfigure}{r}{0.25\textwidth}
    \vspace{-#2}
    \begin{center}
      \includegraphics[width=0.25\textwidth]{img/#1.pdf}
    \end{center}
    \vspace{-#2}
  \end{wrapfigure}
}

% Coloca uma figura menor (sem ser quadrinhos) apertada verticalmente
\newcommand{\figurapequenainlineapertada}[1] {
  \begin{wrapfigure}{r}{0.25\textwidth}
    \vspace{-40pt}
    \begin{center}
      \includegraphics[width=0.25\textwidth]{img/#1.pdf}
    \end{center}
    \vspace{-40pt}
  \end{wrapfigure}
}

% Coloca quadrinhos
\newcommand{\quadrinhos}[1] {
    \figuragrande{quad#1}
}

% Coloca um mapa
\newcommand{\mapa}[1] {
    \begin{figure}[H]
        \centering
        \includegraphics[height=\textwidth, angle=90]{img/#1.pdf}
    \end{figure}
}

%FIXME GAMBIARRA
% Coloca um mapa do IME virando ele do jeito certo
\newcommand{\mapaime}[1] {
    \begin{figure}[H]
        \centering
        \includegraphics[width=\textwidth, angle=180]{img/#1.pdf}
    \end{figure}
}



% ============================ Documento ===============================
\begin{document}

% Capa -------------------------------------------------------------------------
\begin{figure}[p]
  %\begin{center}
    %\hspace{1.7cm}
    \includegraphics[height=1.05\textheight]{img/capa_2019.pdf} %REFTIME
  %\end{center}
\end{figure}
\thispagestyle{empty} % Some com aquele número 1 feio na capa.
\clearpage
\newpage

% Começa a numeração aqui para que o LaTeX entenda qual
% página é ímpar e qual é par antes do índice.
\pagenumbering{arabic}

% Pré-Índice -------------------------------------------------------------------
\begin{editorial}{Editorial}

Bixes, foi difícil chegar até aqui. Vocês
estão meio ou completamente perdidos. Temos apenas uma sugestão: aproveitem essa
etapa.  Façam da sua estadia na USP a melhor época das suas vidas. Vocês verão
que a USP tem muitas e muitas coisas a oferecer. Não se preocupem apenas em
estudar e passar de ano, como vocês fizeram durante sua vida inteira; aproveitem
TUDO (vocês ainda vão descobrir a definição de TUDO). Vocês podem não acreditar
nisso agora, mas saibam que viverão momentos inesquecíveis aqui no IME: alguns
fantásticos, outros deploráveis.

Este guia foi feito para que vocês (bixes que não gostam de ler, apenas olhem as figuras)
possam aprender um pouquinho do que é a USP,
o IME e a vida de universitário que se inicia agora. Ele foi escrito numa forma
descontraída e fácil para que vocês consigam entender; mesmo assim, se pintar
alguma dúvida, vocês podem se dirigir a qualquer veterane, e sua dúvida será
sanada (e, quem sabe, talvez você também comece uma nova amizade).

Note que, ao invés de usarmos as palavras "bixos" ou "bixetes", vamos nos referir
a vocês como "bixes", pois estamos dirigindo palavra a todas as identidades de 
gênero! Por motivos de vocabulário e acessibilidade, usamos a linguagem neutra 
apenas em "bixes" e "veteranes".

%Empurra pro fim da página (altura)
%\vspace{\stretch{1}}

%Linha divisória
\rule{\textwidth}{0.5ex}\rule{2ex}{0.5ex}

%REFTIME
{\large\bf Guia do bixe 2019} \\
Uma publicação da Comissão de Recepção

\paragraph{}
\begin{wrapfigure}{R}{0.25\textwidth}
%REFTIME
% Link para "https://ime.usp.br/~kazuo/guia"
% Esse link contém um contador de acessos que redireciona para o guia correto
% Em 2019: https://ime.usp.br/~kazuo/guiadobixo.pdf
  \begin{center}
    \includegraphics[width=0.22\textwidth]{img/qrcode.pdf}
  \end{center}
  \vspace{-20pt}
  \caption{Baixe aqui o Guia em formato digital!}
  \vspace{80pt}
\end{wrapfigure} 
Muitas pessoas dedicaram seu tempo (e suas férias) para que esse Guia ficasse
pronto: desde o Donald Knuth em 1978, na criação do \TeX\makebox{}, até o
pessoal da Gráfica na correria de ontem à noite (afinal, mais uma vez
entregamos os arquivos atrasados para eles). Esperamos que vocês gostem do Guia
e adotem-no como livro de cabeceira. Por fim, nos avisem de qualquer informação
incorreta ou desatualizada, afinal vocês também são responsáveis por tudo que o
IME oferece a partir de agora.

Lembre: este é o primeiro e último ano de vocês como bixes. Aproveitem!
%\vspace{2em}

%REFTIME (Capa muda todo ano)
%Capa de 2019: ???
%Capa de 2018: ???
%Capa de 2017: ???
%Capa de 2016: public domain
%Capa de 2015:
%{\bf Imagem da Capa:} Viktoria Ridzel -- \href{www.viria13.deviantart.com}{viria13.deviantart.com}

Este trabalho está licenciado sob a Licença Attribution-ShareAlike 4.0 da
Creative Commons. Para ver uma cópia desta licença, visite
\url{https://creativecommons.org/licenses/by-sa/4.0/deed.pt_BR} ou envie
uma carta para Creative Commons, 444 Castro Street, Suite 900, Mountain View,
California, 94041, USA.
\\
\begin{figure}[H]
    \centering
    \includegraphics{img/cc/by-sa.png},
\end{figure}

\end{editorial}


% Carta Aos Ingressantes -------------------------------------------------------
\begin{editorial}{Carta Aos Ingressantes}

Agora que vocês entraram na USP, bixes, vocês adquiriram novas
responsabilidades.  Vocês são responsáveis por si mesmos, isto é, ninguém irá se
preocupar com seus problemas acadêmicos (matrículas, notas erradas, dificuldade
com algumas matérias, rixa com professores etc.) se vocês mesmos não se
preocuparem. Existem pessoas que poderão ajudá-los, mas só o farão se vocês
forem procurá-las. Caso contrário, os únicos prejudicados serão vocês.

A faculdade não é o Paraíso (essa estação fica na Avenida Paulista), mas pode
melhorar a cada dia. Nós, alunos, também devemos contribuir para essa melhora.
Vocês são o futuro da universidade. Portanto, participem, reclamem, busquem seus
direitos, ajudem e, principalmente, não tenham medo de cara feia, pois isso é o
que não vai faltar.

Lembrem-se que vocês não são mais crianças e já sabem o que querem sem que
outros precisem decidir por você; então procurem o que lhes interessa: iniciação
científica, estágios, monitorias, matérias que não são obrigatórias mas que
vocês gostariam de fazer (mesmo que não tenha nada a ver com seu curso);
participação no CAMat ou na Atlética, esportes no CEPE, artigos no jornal, etc.,
etc., etc.

Não é só porque vocês podem, que vocês devem fazer tudo sozinhos, portanto lembrem que
seus amigos vão ser muito importantes para vocês e para o bom andamento do seu
curso. Procurem combinar atividades fora da faculdade, diferentes do cotidiano,
porque isso ajuda a amenizar o estresse que o dia a dia na faculdade pode
trazer.

Além de tudo, vocês não podem se esquecer de uma parte importante de suas vidas
universitárias, que é estudar. Tentem não deixar para estudar na véspera da
prova porque a probabilidade de vocês não irem bem é bem alta (exceto se vocês
tiverem uma sala do tempo em casa, que nem a do Dragon Ball). A mesma coisa se
aplica aos EPs, ainda mais que, quanto mais próximo da data de entrega, mais
erros vão aparecer e, na maioria dos casos, os EPs não são aceitos depois da
data limite de entrega e, caso sejam aceitos, não valerão a mesma nota.

Além disso, como em todo lugar, existem aquelas pessoas ranzinzas e
pentelhas que acham super bacana acabar com a graça de todo mundo,
criticar o IME e aumentar a sua baixa autoestima dizendo que os cursos
são impossíveis e que vocês não vão se formar nunca, mas não acreditem
nelas, vocês entraram aqui com um propósito. Sigam-no.  

Obs.: não se assustem com palavras e siglas que vocês não entenderam ou 
não entenderem! Continuem lendo, porque tudo será explicado em detalhes, 
mastigado, tim-tim por tim-tim. Saibam que vai faltar um monte de siglas...  
Aprendam-nas seletiva e rapidamente. Destaques para: USP, IME, MAC,
MAT, MAE, MAP, BCC, BM, BMA, LIC, BMAC, BE, CEPE, CEAGESP, CTA, RD, CAMat,
AAAMat, SSG, P1, P2, P3, P4, P5, Pn, PQP, CEC, CNPq (=\$), FAPESP
(=\$\$\$), CG, DP, REC, SUB (esta última, ou talvez as duas ou três últimas,
vocês vão conhecer bem melhor, mais cedo ou mais tarde).

\end{editorial}


% Os Dez Mandamentos -----------------------------------------------------------
\begin{editorial}{Os Dez Mandamentos}

\begin{enumerate}
  \item Os veteranes têm sempre razão, exceto quando os bixes têm;
  \item Na hipótese de o bixe ter razão, entra imediatamente em vigor
        o primeiro mandamento;
  \item Em qualquer evento social, as despesas correm sempre por conta dos
        bixes, menos o que os veteranes consomem;
  \item Os bixes têm o direito de permanecerem calados (exceto quando
        interpelados por veteranes ou quando quiserem falar). Tudo o que eles
        disserem pode ser usado contra eles, \textit{c’est la vie};
  \item Os bixes devem se apresentar imediatamente em caso de convocação por
        veteranes, a não ser que estejam fazendo algo melhor da vida. Os bixes
        que não se apresentarem se arrependerão amargamente por perderem
        conselhos valiosos; %não queremos ameaçar ninguém :)
  \item Não são garantidos no IME os direitos constitucionais dos bixes à vida
        social, liberdade de sono e igualdade de notas;
  \item Os bixes devem estar prontos para assumirem as seguintes funções para
        veteranes: completar mesa de jogo, companhia de bandejão, completar
        time(s) etc., mas somente quando tiverem vontade;
  \item Os bixes devem amar e respeitar seus veteranes, a vida, as plantas,
        todo mundo e acima de qualquer coisa, Goku;
  \item Para os casos não abrangidos por estas regras, a decisão final correrá
        por conta dos veteranes, com exceção do que não lhes diz respeito.
  \item Todos os bixes estão despreparados, exceto os que leram isto.
\end{enumerate}

Como bixes, vocês têm todo o direito de reclamar dos mandamentos! Qualquer
reclamação deverá ser protocolada em três vias datadas, assinadas e
autenticadas, com firma reconhecida em cartório, e assim encaminhadas à
Comissão de Recepção 2019 %REFTIME
via mala direta e serão imediatamente incineradas. Alguns veteranes sugeriram
que incinerássemos os reclamantes também. Ao invés disso, cogitamos incinerar
tais veteranes. A medida passou por estudo e, devido aos custos operacionais e
ao lixo tóxico que seria produzido, foi rejeitada.

OU vocês podem conversar diretamente com um veterane da Comissão :)

\end{editorial}


% Os Sete Pecados --------------------------------------------------------------
\begin{editorial}{Os 7 pecados cometidos por bixes}

\begin{itemize}
  \item Pegar o circular errado;
  \item Ir bandejar sem crédito no cartão USP;
  \item Jogar o talher do bandejão no lixo;
  \item Dizer que cursa ``matemática'' (bixes, vocês cursam
        lic/pura/aplicada/bcc/estat/bmac...);
  \item Procurar pelo Jupiter no Google e encontrar uma imagem do planeta;
  \item Esquecer o dia da prova;
  \item Confundir o bilhete único com o BUSP e pagar o circular.
  \item Comprar na lanchonete;
  \item Esquecer de fazer a matrícula;
  \item Atrasar a devolução de livros na biblioteca;
  \item Não ter moeda para a \sout{máquina de café}\footnote{R.I.P. nossa querida Jennifer,
    conheça-a em \url{fb.com/jennifer.maccafe}};
  \item Entregar a versão errada do EP, ou pior, entregar em \texttt{*.doc};
  \item Não saber que o professor cancelou a aula;
  \item Achar que o P3 é perto;
  \item Frequentar a turma errada da matéria;
  \item Ir para a aula só assinar a lista, e o professor não passar lista;
  \item Desperdiçar quota na impressora sem papel;
  \item Deixar o computador logado na Linux ou no CEC e ser trollado;
  \item Acender a luz da vivência antes das 8h da manhã;
  \item Atrapalhar a aula gritando alto demais no truco;
  \item Se matricular em uma matéria em um campus do interior sem querer;
  \item Esquecer do prazo de trancamento;
  \item Não saber quem é o Goku;
  \item Não saber contar quantos pecados tem na lista de 7 pecados.
\end{itemize}

\end{editorial}



%Renomeando o índice------------------------------------------------------------
\renewcommand{\contentsname}{\center Esse guia contém...}

% Coloca o índice --------------------------------------------------------------
\tableofcontents
% Remove números de página do índice
\thispagestyle{empty}
\addtocontents{toc}{\protect\thispagestyle{empty}}
\addtocontents{toc}{\protect\pagestyle{empty}}
\newpage
\thispagestyle{empty}
\cleardoublepage

% Começa a numeração aqui
\pagenumbering{arabic}


% Comissão de Trote e o Kit-bixo -----------------------------------------------
\begin{secao}{Comissão de Recepção e \textit{Kit}-bixe}

A Comissão de Recepção aos Calouros, também conhecida apenas como Comissão de 
Recepção, é responsável por auxiliar os ingressantes em seus primeiros momentos 
IMEanos. Sabemos que não é um momento fácil. Vocês estão entrando em uma nova fase 
de suas vidas, em um lugar estranho, com pessoas estranhas (em todos os sentidos),
e nosso objetivo é fazer com que vocês se sintam bem-vindos e se integrem
($\int$) com seus coleguinhas e seus veteranes!

A Comissão de Recepção organiza a super Semana de Recepção, cheia de atividades
legais! (Vocês devem ter em mãos a programação da Semana de Recepção. 
Se não receberam, procurem arrumar uma, logo!).
Somos nós também que organizamos o magnífico encontro dos bixes: AnIME-se, que 
deverá ocorrer no dia \textbf{16 de Março}. Não esqueçam essa data, bixes! %REFTIME

A Comissão de Recepção é formada pelos mais animados e divertidos veteranes. Como
dissemos, eles organizam a Semana de Recepção, ajudam bixes perdidos, mantêm um clima
alegre durante recepção, fazem isso e aquilo...
Vocês devem estar pensando ``Puxa! Como nós, bixes, podemos retribuir
tamanha dedicação?'' É simples, bixes: {\bf\em comprem o \textit{kit}-bixe}!!!

O \textit{kit}-bixe, como vocês devem saber, é um conjunto de coisas
importantíssimas para vocês, ingressantes perdidos! Ele contém dois tipos de
itens:
\begin{itemize}
\item itens úteis;
\item itens essenciais.
\end{itemize} %REFTIME (Isso muda todo ano)
Dentre eles, temos uma camiseta com uma linda estampa do tema do ano de vocês,
que serve para que nós os identifiquemos como bixes e para que as pessoas na 
rua achem que vocês são inteligentes; materiais personalizados com e sem o símbolo do IME para 
vocês usarem nas aulas e mostrarem aos seus amigos; uma caneca (também personalizada) 
para vocês economizarem muitos copos no bandejão; um chaveiro personalizado muito 
maneiro que também serve para abrir garrafas;
dois adesivos, um para vocês colarem no carro que vão ganhar de presente, e outro
para os papais (quando vocês virem o adesivo vocês vão entender); calculadora para 
usarem nas suas provas de Estatística e um caderno super
lindo e sem pauta para vocês desenharem seus gráficos das aulas de Cálculo. Tudo
isso preparado com muito carinho e especialmente para vocês pela Comissão de Recepção,
além de estarem contidos em uma mochila sport do IME-USP, afinal, não dá pra
carregar tudo isso na mão, né?

Adquiram o maravilhoso \textit{kit}-bixe do IME-USP. Ele estará à venda na Semana de Recepção
e inclusive na semana subsequente, mas somente enquanto durarem os estoques. (Não vão chorar
depois, hein?)

Para terminar, além de tudo isso, a Comissão de Recepção é que faz esse maravilhoso guia que
vocês estão lendo agora (ou estão só olhando as figuras, vai
saber...). Esperamos que vocês gostem do nosso trabalho! Qualquer coisa, nos
procurem! Entrem na nossa página no Facebook: {\tt fb.com/recepcaoimeusp} contando
o que vocês sentiram ao ler o guia, o que fizeram de errado com vocês na
recepção (com o nome do veterane na denúncia), o que vocês acharam da Semana
de Recepção e se vocês se sentiram bem-vindos ou querem voltar logo para perto
da sua mãe. Estaremos sempre prontos a ajudá-los!
\end{secao}

\quadrinhos{1}

% O CAMat ----------------------------------------------------------------------
\begin{secao}{O CAMat}

O Centro Acadêmico da Matemática, Estatística e Computação é o que chamamos de
CAMat. É um órgão reconhecido pelo IME, feito pelos estudantes e para os
estudantes.

É dever do CAMat lutar pelos nossos direitos, organizando debates, palestras,
além, é claro, de eventos culturais, encontros, feiras de livros, festas...

O CAMat representa os alunos junto ao IME e todas as demais entidades do nosso
Instituto. O CAMat não é simplesmente uma sala, um espaço. O CAMat somos nós,
todos os alunos do IME, anualmente representados por um grupo de alunos que
vence um pleito democrático, do qual você também participará, seja votando ou,
logo, logo, sendo votado!

Todas as atividades promovidas pelo CAMat são discutidas e decididas em
reuniões abertas, nas quais, pelo menos desde a gestão 2010, todo aluno do IME tem
direito a voz e voto. Elas serão a melhor oportunidade para você apresentar suas
ideias, questionar, ou simplesmente ficar a par do que o CAMat
anda fazendo.

A sala do CAMat  fica dentro da Vivência, na sala 18 do bloco B. Você tem acesso
livre lá dentro para conversar, ver TV, ler jornais e revistas, jogar sinuca,
pebolim, xadrez, videogame, BARALHO (aqui você vai aprender um número inimaginável de
jogos de baralho...)\footnote{Veja a seção desse guia dos jogos da
Vivência} e até estudar.

Para salvar suas costas dos milhares de livros de defesa contra as artes das
Trevas (Cálculo, Álgebra, Análise...) que você precisará carregar, o CAMat
também aluga armários aos alunos a uma taxa insignificante. Fiquem atentos, pois no
início das aulas do 1º semestre ocorrerá a renovação dos armários com
aqueles alunos veteranes que já os têm e certamente sobrará uma boa quantidade
de armários para outros, incluindo vocês, bixes.

Além desses, o CAMat oferece alguns outros serviços:

\begin{itemize}
  \item Organizamos um Banco de Provas, que contém várias provas dos anos
    anteriores feitas pelos maravilhosos veteranes. Não se esqueça de contribuir
    com sua prova também.
  \item Emprestamos baralhos, violões e fichas de pôquer mediante apresentação
    de carteirinha USP.
  \item Em seus momentos de fome e sede fique tranquilo, pois no CAMat você
    também encontra alguns "comes e bebes", prontos para matar aquela fome e
    sede que teimam em aparecer praticamente a toda hora... Ah! Os preços são
    praticamente de atacado!
  \item Temos também uma máquina de café disponível a todos que queiram lutar
    contra o sono... É só chegar lá e preparar, de graça!
  \item Plantões de dúvidas para dar aquela força nas disciplinas! (veja no
    final desta seção)
\end{itemize}

Bom, mas o CAMat não é e nem deve ser feito só de serviços. Há atualmente
muitas mudanças no nosso instituto, e precisamos garantir que os alunos não
saiam prejudicados. Em 2010, trouxemos, junto com o
DCE, uma exposição sobre a ditadura militar para o saguão do IME. Para abrir a
exposição e discutir o assunto com a gente, recebemos o então Ministro dos
Direitos Humanos, Paulo Vannuchi. Achamos fundamental discutir assuntos não só
relativos ao IME, mas também à USP, ao Brasil e ao mundo.

O CAMat também tem várias formas de você entrar em contato sempre que precisar
(quando roubarem seu lanche, puxarem seu cabelo ou te chamarem de bobo):

\begin{description}
\item [Página:] \url{http://www.ime.usp.br/~camat}
\item [e-mail:] camat.usp@gmail.com
\item [grupo de e-mails:] camat-aberto@googlegroups.com.br (faça parte!)
\item [Facebook:] \url{fb.com/CAMatUSP}
\item [Telefone:] 3091-6293
\end{description}

LEMBRE: não deixem de participar (ao menos para conhecer) das reuniões do
CAMat. Local e horários serão propriamente divulgados.

LEMBRE$^2$: vocês são SEMPRE bem-vindos na salinha do CAMat.

\begin{subsecao}{Plantão de dúvidas}

Pra quem tá perdido no curso, sentindo falta de alguma coisa que não aprendeu
direito no colégio, não está conseguindo acompanhar as disciplinas e
eventualmente quer um conselho de alguém que está há mais tempo no IME, temos o
Plantão de Dúvidas!

Ao longo do semestre estudantes do próprio IME (e às vezes professores) se
voluntariam pra dar aulas com assuntos de ensino médio com que vocês podem não
estar confortáveis e são importantes: funções, trigonometria, logaritmos etc.

O próprio instituto, pela Comissão de Graduação, nos ajuda todos os anos a
encontrar salas e voluntários pra essas aulas, mas os detalhes de dias e
horários são decididos apenas quando começa o semestre, de forma que não temos
como já informar onde serão esses plantões.

Portanto, fiquem de olho nos seus e-mails institucionais, murais e na seguinte
página: \url{https://www.facebook.com/plantaoduvidas/}

Não hesite em nos procurar!

\end{subsecao}

\end{secao}

\quadrinhos{2}

% A Atlética -------------------------------------------------------------------
\begin{secao}{A Atlética} %REFTIME - Quase tudo aqui é REFTIME...

\begin{subsecao}{O que é a AAAMat?}

É a Associação Atlética Acadêmica da Matemática - a entidade mais divertida do
mundo!!!! - que tem como objetivo trazer os melhores momentos da sua vida
universitária! A Atlética é formada por um grupo de IMEanos (gestão) que é
responsável por organizar atividades esportivas e eventos (festas, pizzadas,
premiações etc) para a comunidade IMEana.

\end{subsecao}
%%%%%%%%%%%%%%%%%%%%%%%%%%%%%%%%%%%%%%%%%%%%%%%%%%%%%%%%%%%%%%%%%%%%%%%%%%%%%%%%

Algumas das atividades da Atlética:

\begin{subsecao}{Atividades esportivas internas:}

Na AAAMat existem diretores de modalidade (DMs) que são pessoas responsáveis
pelos treinos e campeonatos dos seguintes esportes: futebol de campo, futsal,
basquete, vôlei, handebol, atletismo, natação, tênis de mesa, tênis de campo,
xadrez, judô, beisebol, softbol, ultimate frisbee, bridge, rugby e e-sports
(\textit{League of Legends}, \textit{Hearthstone} e
\textit{Counter Strike: Global Offense}). Além disso, contamos com a nossa 
amada BatIMEduca, bateria do IME juntamente com a Pedago.

Os dias e horários dos treinos/jogos de cada uma dessas modalidades serão
sempre informados através do site e do mural da Atlética, localizado na entrada
do bloco B (aquele mural verde, colado na parede em frente a lanchonete).

Com essa variedade de modalidades, não tem desculpa pro sedentarismo hein,
bixe? Se você não conhece nenhuma delas, a gente te apresenta e se já conhece,
vem dar aquele “ooooi, sumido!!!” pra aquele esporte que você largou por conta
da fuvest! Ah, não esqueça de torcer pelos nossos atletas! Raça e coração, IME!

Além disso, a Atlética promove anualmente campeonatos internos daqueles jogos
que a gente passa hooooras fritando em casa. Já foram promovidos campeonatos de
Winning Eleven, LoL, Mario Kart, Super Smash Bros, Mario Tenis, Guitar Hero e
sinuca.

Ideias e sugestões sobre novas modalidades, campeonatos, inters etc. são
sempre muito bem-vindas! Conversem com a gente!

\end{subsecao}
%%%%%%%%%%%%%%%%%%%%%%%%%%%%%%%%%%%%%%%%%%%%%%%%%%%%%%%%%%%%%%%%%%%%%%%%%%%%%%%%

A AAAMat também representa o IME em diversos campeonatos universitários. São
eles:

\begin{subsecao}{BichUSP}

De nome intuitivo e charmoso, o BichUSP é um campeonato disputado entre as
faculdades da USP em que apenas os bixes (VOCÊS!) participam. O campeonato
acontece logo nas primeiras semanas de aula, sempre aos finais de semana. Aqui,
vocês têm a chance de suar a camisa IMEana pela primeira vez e ver a torcida
indo ao delírio em cada jogada - ganhando ou perdendo, seus veteranes estarão
vibrando por vocês!

%REFTIME
Esse ano o BichUSP acontecerá nas seguintes datas:

\begin{itemize}
  \item Tênis todos os dias do BichUSP
  \item 16 e 17/03/2019 - Basquete e Handebol.
  \item 23 e 24/03/2019 - Natação, Atletismo, Tênis de Mesa e Xadrez.
  \item 30/03 e 01/04/2019 - Futebol de campo e Rugby.
  \item 06 e 07/04/2019 - Vôlei e Futsal.  
\end{itemize}

``Mas, Atlética, eu não sei jogar nenhum desses esportes :('' - Não tem
problema, a gente te ensina! Teremos treinos especiais para que vocês conheçam
a modalidade, os DM’s, os técnicos, a gente e os outros bixes que te
acompanharão nesse momento único da graduação!

Se você acha que tem alergia a esportes, dá uma chance da gente te mostrar o
contrário! São várias modalidades com várias dinâmicas diferentes, alguma delas
com certeza vai se encaixar no que você gosta! Se quiser só assistir no começo
e vir torcer com a gente, apareçam nos jogos que nós gritamos "VERMELHO E
BRANCO ATÉ MORRER" todos juntos!

Pra vocês se inspirarem, fizemos essa tabelinha que mostra quantos bixes
brilharam em anos anteriores. Estamos ansiosos pra completar ela com as
conquistas que virão esse ano:

%REFTIME
\begin{center}
  \begin{tabular}{c|c}
    \hline
    Ano & Campeão\\
    \hline
    2005 & Basquete Masculino \\
    2005 & Tênis de Mesa Feminino \\
    2007 & Atletismo Masculino\\
    2009 & Atletismo\\
    2011 & Tênis de Campo Feminino\\
    2012 & Basquete Feminino\\
    2014 & Tênis de Mesa Masculino\\
    2014 & Futsal Masculino\\
    2015 & Futsal Masculino\\
    2016 & Vôlei Masculino\\
    2017 & Xadrez, Rugby Misto e Rugby Feminino (IME+EEFE)\\
    2018 & Xadrez, Futebol de Campo F.\\
    2019 & VEM BIXES!\\
    \hline
  \end{tabular}
\end{center}

\end{subsecao}
%%%%%%%%%%%%%%%%%%%%%%%%%%%%%%%%%%%%%%%%%%%%%%%%%%%%%%%%%%%%%%%%%%%%%%%%%%%%%%%%
\begin{subsecao}{Copa USP}

A Copa USP é o primeiro campeonato após o BichUSP, e existem duas séries (Azul:
1ª divisão e Laranja: 2ª divisão). São nesses jogos que colocamos em prática
tudo o que fizemos nos treinos semanais para brilharmos nos jogos da fase de
grupos e então seguir arrasando nos jogos mata-matas.

\end{subsecao}
%%%%%%%%%%%%%%%%%%%%%%%%%%%%%%%%%%%%%%%%%%%%%%%%%%%%%%%%%%%%%%%%%%%%%%%%%%%%%%%%
\begin{subsecao}{Jogos da Liga}

Acontece no segundo semestre (UFA, já passou a P1 de cálculo, vem Cálculo 2!!!
\sout{Ou não}) e nessa competição não existe separação por séries. Somente as
faculdades da USP jogam, as disputas são sorteadas para formarem grupos e
então apenas os melhores colocados seguem para a fase final. Essa é a
oportunidade perfeita para gritar um "CHUPA POLI" na arquibancada.

\end{subsecao}
%%%%%%%%%%%%%%%%%%%%%%%%%%%%%%%%%%%%%%%%%%%%%%%%%%%%%%%%%%%%%%%%%%%%%%%%%%%%%%%%
\begin{subsecao}{NDU}

Esse campeonato acontece duas vezes ao ano, e várias faculdades de São Paulo
(tanto das USP quanto algumas não-USP) competem na cidade em busca dos melhores
resultados. Confira com o DM da modalidade se o time está participando da
competição.

\end{subsecao}
%%%%%%%%%%%%%%%%%%%%%%%%%%%%%%%%%%%%%%%%%%%%%%%%%%%%%%%%%%%%%%%%%%%%%%%%%%%%%%%%
\begin{subsecao}{BIFE}

O BIFE É O MELHOR EVENTO ESPORTIVO DO MUNDO! As iniciais das quatro fundadoras
(Bio, IME, Fau e Eca) formam a sigla que dá nome a esses jogos universitários
que a gente tanto ama. Trata-se de um campeonato entre nove faculdades da USP:
a VET, GEO, Física, FFLCH, Química e, é claro, as quatro fundadoras já citadas.

Funciona assim: em um determinado feriado, jogadores, torcedores, festeiros e
simpatizantes se deslocam até alguma cidade do interior do Estado. A cidade nos
dá um alojamento (leia-se: local para tomar um banho quentinho e descansar no
aconchego de sua barraca), alguns ginásios e um local para as festas. São
quatro dias muito divertidos e engraçados, onde há rivalidade apenas dentro de
quadra - porque fora é muito amor e integração!

Nosso histórico neste Inter é de parar o trânsito! Olhem só:

%REFTIME
\begin{center}
  \begin{tabular}{c|c|c}
   Ano & Cidade & Campeão\\
   \hline
   1999 & Jacareí & IME\\
   2000 & Não Houve & - \\
   2001 & Serra Negra & IME\\
   2002 & Socorro & ECA\\
   2003 & São Sebastião & IME\\
   2004 & Cruzeiro & FFLCH\\
   2005 & Jacareí & FFLCH\\
   2006 & Lorena & IME\\
   2007 & Piedade & IME\\
   2008 & Itapeva & IME\\
   2009 & Cruzeiro & IME\\
   2010 & Barra Bonita & IME\\
   2011 & Casa Branca & IME\\
   2012 & Barra Bonita & IME\\
   2013 & Sumaré & ECA\\
   2014 & Cidade/Araraquara & FFLCH\\
   2015 & Taquaritinga & FFLCH\\
   2016 & Registro & FFLCH\\
   2017 & Avaré & FFLCH\\
   2018 & Casa Branca & ICBIÓ\\
   2019 & ??? & VAMO IME!!!
  \end{tabular}
\end{center}

%REFTIME
Nesses últimos seis anos não conseguimos o título MAS ESSE ANO VAI! Nossos 
times contam com vocês para que juntos possamos ser hendecacampeões! (essa
palavra existe mesmo)

\end{subsecao}
%%%%%%%%%%%%%%%%%%%%%%%%%%%%%%%%%%%%%%%%%%%%%%%%%%%%%%%%%%%%%%%%%%%%%%%%%%%%%%%%
\begin{subsecao}{Títulos}

Fruto de muito treino, empenho, suor, torcida e amor pelo IME-USP, reunimos
abaixo algumas de nossas conquistas:

%REFTIME
\begin{center}
  \begin{tabular}{c|c|c|c}
    Ano & Campeonato & Modalidade & Colocação\\
    \hline
    2006 & Jogos da liga  & Handebol Masc.  & 2º\\
    2009 & Copa USP       & Handebol Masc.  & 1º\\
    2012 & Copa USP       & Handebol Masc.  & 2º\\
    2016 & IMEACHECA      & Handebol Masc.  & 2º\\
    2018 & Jogos da Liga  & Handebol Masc.  & 2º\\
    2017 & Copa USP       & Baseball Fem.   & 2º\\
    2011 & BOBPAI         & Baseball        & 1º\\
    2016 & BOBPAI         & Baseball        & 2º\\
    2016 & Liga Paulista  & Baseball        & 3º\\
    2016 & Wakaba         & Softball        & 2º\\
    2016 & Softparty      & Softball        & 1º\\
    2011 & Jogos da liga  & Basquete Fem.   & 1º\\
    2012 & Copa Camp      & Basquete Fem.   & 2º\\
    2012 & Copa USP       & Basquete Fem.   & 3º\\
    2014 & Interfarofa    & Futsal Fem.     & 1º\\
    2015 & NDU            & Futsal Fem.     & 2º\\
  \end{tabular}
\end{center}
\begin{center}
  \begin{tabular}{c|c|c|c}
    Ano & Campeonato & Modalidade & Colocação\\
    \hline
    2017 & Copa USP       & Futsal Fem.     & 1º\\
    2017 & Jogos da Liga  & Futsal Fem.     & 2º\\
    2017 & IMEACHCA       & Futsal Fem.     & 1º\\
    2015 & Camp. G-4      & Vôlei Masc.     & 2º\\
    2016 & Copa USP       & Vôlei Masc.     & 1º\\
    2017 & Copa USP       & Vôlei Masc.     & 2º\\
    2017 & NDU            & Vôlei Masc.     & 1º\\
    2018 & Copa USP       & Vôlei Masc.     & 1º\\
    2018 & BIFE           & Vôlei Masc.     & 2º\\
    2018 & NDU            & Vôlei Masc.    & 2º\\
    2015 & Camp. G-4      & Handebol Fem.   & 1º\\
    2015 & Integramix     & Handebol Fem.   & 1º\\
    2016 & CUPA           & Handebol Fem.   & 2º\\
    2016 & Interfarofa    & Handebol Fem    & 1º\\
    2018 & Copa USP       & Handebol Fem.   & 1º\\
    2009 & Intercalouros  & Atletismo       & 1º\\
    2011 & LUPAA          & Atletismo       & 1º\\
    2015 & Integramix     & Futebol Campo   & 1º\\
    2017 & Copa USP       & Futebol Campo M & 1º\\
    2011 & Copa USP       & Futsal Masc.    & 1º\\
    2012 & Copa Camp      & Futsal Masc.    & 1º\\
    2013 & NDU            & Futsal Masc.    & 1º\\
    2013 & Jogos da Liga  & Futsal Masc.    & 2º\\
    2015 & Integramix     & Futsal Masc.    & 1º\\
    2015 & NDU            & Futsal Masc.    & 2º\\
    2016 & IMEACHECA      & Vôlei Fem.      & 1º\\
    2016 & Gran Prix USP  & Vôlei Fem.      & 3º\\
    2017 & Gran Prix USP  & Vôlei Fem.      & 3º\\
    2017 & Copa USP       & Jiu-jitsu       & 2º\\
    2017 & Jogos da Liga  & Tênis Campo M   & 2º\\
    2017 & NDU            & Xadrez          & 1º\\
    2017 & Copa USP       & Xadrez          & 3º\\
    2017 & Jogos da Liga  & Xadrez          & 2º\\
    2017 & TUES           & Hearthstone     & 2º\\
    2018 & BIFE           & Basquete Masc.  & 1º\\
    2018 & Jogos da Liga  & Basquete Masc.  & 2º\\
    2018 & BIFE           & Natação Masc.   & 2º\\
    2018 & Jogos da Liga  & Natação Masc.   & 2º\\
    2018 & BIFE           & Rugby Masc.     & 2º\\
    2018 & BIFE           & Futebol Campo F.& 2º
  \end{tabular}
\end{center}

\end{subsecao}
%%%%%%%%%%%%%%%%%%%%%%%%%%%%%%%%%%%%%%%%%%%%%%%%%%%%%%%%%%%%%%%%%%%%%%%%%%%%%%%%
\begin{subsecao}{Outras atividades}

\begin{subsubsecao}{Vendas}

A Atlética também quer te ajudar a vestir o vermelho e branco (que, a essa
altura, já corre em suas veias! :D) e traz pra você diversos produtos
personalizados, tais como: adesivos, tatuagens, canecas, talabartes, chaveiros,
samba-canção, cadernos, estojos, mousepads, chinelos, camisetas e agasalhos do
IME, pra você sair por aí esbanjando seu amor pelo IME-USP $<$3

Chegou na hora da prova e esqueceu seu kit bixe em casa? A gente te salva --
também vendemos lápis, borracha, régua, caneta e calculadora!

\end{subsubsecao}
%%%%%%%%%%%%%%%%%%%%%%%%%%%%%%%%%%%%%%%%%%%%%%%%%%%%%%%%%%%%%%%%%%%%%%%%%%%%%%%%
\begin{subsubsecao}{Festas}

%REFTIME
A Atlética e o CAMat já promoveram muitas festas e happy hours. Atualmente,
promovemos a I Will SurvIME, a FofIMEduca (Com a galera da Pedagogia), o 
JunIME, a Melhores do Ano, o sarau e alguns HHs durante o ano e auxiliamos 
as festas pré-BIFE (Desmame, Engorda e Abate). Todas imperdíveis!
Esperamos vocês!

\end{subsubsecao}
%%%%%%%%%%%%%%%%%%%%%%%%%%%%%%%%%%%%%%%%%%%%%%%%%%%%%%%%%%%%%%%%%%%%%%%%%%%%%%%%
\begin{subsubsecao}{Como falar com a Atlética?}

A salinha da AAAMat é a B-18. Ela fica dentro da vivência e é pequena, mas
sempre cabe mais um! Sempre que precisarem conversar com a Atlética, vocês
podem ir até lá e falar com qualquer membro da gestão. Além disso, também temos
outros meios de contato, tais como o telefone: (11) 3091-6378 ou o e-mail:
atletica@ime.usp.br

Se quiserem ficar por dentro de tudo que acontece na atlética, vocês também
podem:

\begin{itemize}
  \item Acompanhar o site da Atlética: \url{www.ime.usp.br/~atletica}
  \item Curtir nossa página no Facebook: \url{fb.com/aaamat.ime}
  \item Seguir a gente no Instagram: @aaamat\_imeusp
  \item Participar da nossa lista de e-mails: basta enviar um e-mail em branco
        para \url{aaamat-diretoria+subscribe@googlegroups.com} e seguir as
        instruções enviadas pro seu e-mail.
\end{itemize}

%REFTIME
A Atlética inicia 2019 esperando vocês, bixes lindos, para que juntos possamos
trazer muitos títulos e troféus para casa! É importantíssimo que vocês saibam
que estamos abertos para qualquer tipo de crítica, dúvida, ideia ou sugestão.

Vocês são SEMPRE muito bem-vindos em nossa sala, atividades, times e eventos!
\end{subsubsecao}
\end{subsecao}
\end{secao}


% Dissecando os Cursos ---------------------------------------------------------
\input{dissecando_os_cursos.tex}
\quadrinhos{6}

% Júpiter Web ------------------------------------------------------------------
\begin{secao}{JupiterWeb}

O \sout{Luciferweb} Jupiterweb é o sistema que administra (quase) tudo que
vocês, bixes, fazem na sua graduação. Desde as suas matrículas até a emissão de
um documento para pagar meia no cinema, o Jupiterweb vai ser muito útil em
vários momentos se vocês conhecerem tudo que ele tem a oferecer. Por isso,
assim que puderem, procurem explorar cada uma de suas funcionalidades.

\begin{subsecao}{Calendário}

Muito útil pra saber quando é \sout{feriado} que abre a matrícula ou o período
de requerimento.

\end{subsecao}

\begin{subsecao}{Cursos de ingresso}

Aqui, você pode ver as informações sobre o seu curso, o projeto pedagógico e
a grade curricular. Para isso, basta inserir a unidade do seu curso (no caso,
IME) e o nome do seu curso.

Preste atenção na grade curricular, pois ela tem todas as disciplinas que
vocês terão que cursar durante a graduação. Elas estão separadas em obrigatórias
(matérias que terão que cursar para se formar), optativas eletivas (matérias
que terão que escolher dentre algumas opções pra fazer) e optativas livres
(aquelas que acham que você pode se interessar por fazer).

Tome cuidado principalmente com as obrigatórias que são pré-requisitos de
alguma outra matéria! Você não poderá cursar uma matéria se não tiver passado
nos seus respectivos pré-requisitos. Então, \textbf{passem em Cálculo!!!!}

Existem também os requisitos fracos (mas são uma raridade). Neles, vocês não 
precisam ter passado na matéria para cursar a outra, basta ter tirado 3.

\end{subsecao}

\begin{subsecao}{Disciplinas}

Mostra todas as disciplinas da USP (até as que não estão mais disponíveis).
Vocês podem pesquisar por código, por algum trecho do nome ou por unidade. 
Além de ser útil para conferir algumas informações sobre aquela matéria que 
você precisa cursar, ainda ajuda a procurar matérias diferentes pela USP 
(como “Pintura e suas técnicas”, “Psicologia da Morte” ou “Pesca Sustentável”).

\end{subsecao}

\begin{subsecao}{Grade Horária}

Lista as disciplinas em que vocês estão inscritos/matriculados/pendentes, com os
horários de cada uma. É onde geralmente vocês conferem se tá tudo certo com a
sua matrícula ou não. Quando as aulas começam, o ideal é que apareça
“Matriculado” (a matéria fica cinza). Se alguma disciplina aparecer como
pendente (fica azul), então corram na seção de alunos para conferir o que está
errado!

\end{subsecao}

\begin{subsecao}{Resumo Escolar}

O resumo escolar é um documento enorme com tudo de bom (e de ruim) que
aconteceu na sua vida acadêmica na USP. Lá vocês podem conferir cada
disciplina cursada (com nota), as disciplinas que vocês trancaram, seus
aproveitamentos de estudo, quantos créditos vocês já fizeram e a tão mística
“Média ponderada”.

\end{subsecao}

\begin{subsecao}{Acompanhamentos}

\begin{itemize}
  \item \textbf{Evolução no Curso}

    É o jeito visual de acompanhar as disciplinas que vocês já fizeram ou que
    estão faltando. Sua principal função é ser printado ao final do curso (com
    tudo verde) para postar no Facebook e fazer sucesso com a família na
    internet.
  \item \textbf{Rendimento Acadêmico}

    Mais um lugar pra vocês conferirem suas notas, quantidade de matérias
    cursadas, e quantidade de créditos cursados, só que aqui tudo aparece em
    forma de gráfico.
  \item \textbf{Perfil de deficiências}

    Questionário da USP sobre deficiências.

  \item \textbf{Meus benefícios e bolsas}

    Onde ficam registradas as suas bolsas e alguns serviços prestados na
    universidade que valem dinheirinhos. Também mostra as datas que a USP lhes
    paga para vocês organizarem sua vida financeira!!

  \item \textbf{Dados pessoais}

    O nome é bem intuitivo. Mantenham tudo atualizado para não ter problemas com
    a SPTrans, SAS ou qualquer coisa que confira seus dados em algum momento.

\end{itemize}

\end{subsecao}

\begin{subsecao}{Emissão de Documentos}

Emite alguns documentos que precisamos no dia a dia, como um comprovante de
matrícula ativa. Dependendo do documento, vocês vão precisar ir até a seção de
alunos, mas os mais requisitados vocês encontram por aqui.

\end{subsecao}

\begin{subsecao}{Matrícula}

Aqui começa uma das partes mais legais da Universidade: A matrícula! Pois é,
bixes, vocês não vão cursar apenas o que a sua instituição de ensino manda, mas o
que vocês quiserem! Apesar de ter uma grade de matérias obrigatórias, vocês
podem cursá-las na ordem que lhes for mais conveniente e sempre ir intercalando
com optativas eletivas e optativas livres.

Não que isso seja uma boa prática. O seu curso foi pensado e montado com uma
sequência de matérias que faça sentido na sua vida, mas vocês têm independência
para quebrar isso caso seja melhor pra vocês.

\textbf{Vocês não precisam se preocupar com isso agora! As matérias do primeiro
semestre já são inseridas automaticamente, vocês só precisam fazer a própria
matrícula a partir do segundo semestre.}

Quando a hora chegar, vocês vão clicar em “Fazer matrícula” e escolher o tipo de
matéria que vocês querem cursar: Obrigatórias, Optativas eletivas, Optativas
Livres ou Extracurriculares. Quando vocês escolhem um desses blocos, vão
aparecer uma série de matérias disponíveis pra vocês com as turmas (cada turma
tem um horário e um professor específico). Basta vocês irem escolhendo quais
querem cursar, conferir umas 15 vezes e salvar (não se esqueçam de salvar,
bixes, ou não estarão matriculados nas disciplinas!).

Pode até parecer um monstro de 7 cabeças, mas está beeeeem longe disso!
Aproveitem esse momento para consultar aquele veterane que você ama muito e que
vai te mostrar como é tudo bem tranquilo.

Mas quando é mesmo que eu faço isso?

O Jupiterweb vai abrir algumas “Interações de matrícula”, geralmente duas ou três.
Apenas nesses dias, o botãozinho de “Fazer matrícula” fica disponível. Vocês
recebem e-mails da USP, fica um banner no meio do saguão do bloco B, fazem
evento no facebook e divulgam a data no calendário da USP… Não tem desculpa pra
não saber quando é!

No dia de abertura das interações, o Jupiterweb revela suas verdadeiras garras
(e lentidões) dando absolutamente todos os erros possíveis e tentando
atrapalhar a sua vida. Então, talvez esse não seja o melhor horário pra tentar
se matricular. A ordem de matrícula não importa, então podem ficar tranquilos e
fazer em qualquer momento de qualquer uma das interações (MAS NÃO SE ESQUEÇAM DE
FAZER EM ALGUM MOMENTO!!!).

Depois de cada interação, o Jupiterweb vai colocar vocês em um ranking que olha seu
semestre, o que vocês deveriam estar cursando, suas notas e o que mais for
conveniente para decidir se vocês ganham ou não uma vaga na disciplina.
Geralmente isso não é um problema com matérias obrigatórias (principalmente se
vocês seguirem sua grade ideal direitinho) ou matérias com poucos inscritos. Mas
se vocês estiverem tentando uma optativa muito concorrida (como ``Psicologia da
morte''), se preparem que vocês podem perder uma vaga muito sonhada.
Esse resultado pode ser acessado no ``Resultado das consolidações''.

\end{subsecao}

\begin{subsecao}{Requerimentos}

Por vários motivos, vocês podem ter sua matrícula de uma disciplina negada ou
até mesmo nem conseguir se matricular porque essa disciplina não tem vagas
reservadas para o seu curso. Isso não quer dizer que vocês não vão conseguir
cursar, mas que vocês vão ter que enfrentar mais burocracias para isso. Os
requerimentos, também conhecidos como melhores amigos do aluno da graduação,
são pedidos formais em que vocês inserem a matéria da qual vocês foram
chutados/negados/impedidos de cursar e uma boa justificativa. Esse requerimento
vai para o coordenador do seu curso e possivelmente para o professor de
matéria, pessoa que vai decidir se o seu motivo é bom o suficiente e se o
número de vagas comporta a sua presença.

Vocês podem fazer um requerimento para cursar qualquer matéria, por qualquer
motivo, mas se o pedido for muito esdrúxulo, vocês podem receber um não.

O resultado dos requerimentos demora 84 anos e um pouco para sair. Por isso,
vocês vão acompanhar as aulas da disciplina requerida mesmo sem saber o
resultado, confiantes de que vocês vão receber um OK depois de um tempo. Existe
um prazo limite para aprovar ou negar um requerimento, então fiquem tranquilos
que até a P1 ou P2, vocês já devem saber se estão matriculados ou não.

Dica: Conversar com os coordenadores e o professor responsável da disciplina
pode ajudar! Se vocês querem muito cursar uma matéria, não tenham vergonha de
lutar por ela :)

\end{subsecao}

\end{secao}


% Um Pouco Sobre o IME ---------------------------------------------------------
\begin{secao}{Um Pouco Sobre o IME}

O IME, Instituto de Matemática e Estatística (não, bixes, não está errado.
Computação não faz parte da sigla mesmo! Aliás, faz sim, mas só os inteligentes
podem ver!), tem vários blocos: bloco A, bloco B, bloco C, bloco D e o
novo bloco, até então um mistério para todos. Leia atentamente, porque com essa
confusão toda, até alguns veteranes insistem em dizer que só existem os blocos
A, B e C.


\begin{subsecao}{Bloco A}
No bloco A, as coisas mais importantes são: as salas dos professores, a parte
administrativa (diretoria, secretarias de departamento), algumas salas de aula
(para a pós-graduação), a IMEjr, a biblioteca, as ET's e a Rede Linux.

\begin{itemize}

\item {\bf IMEjr:} é a empresa júnior do IME, que é administrada pelos próprios
alunos da graduação. Fica na sala 258A.

\item {\bf Biblioteca:} fica logo na entrada, do lado direito. Há algumas mesas
individuais e sofás (ambos muito confortáveis para tirar um cochilo) e algumas salas
com lousa para estudo em grupo. Tem também uma sala de estudos enorme, 
mas lá não pode fazer barulho (cuidado até mesmo ao abrir e fechar a porta!).
Veja mais informações sobre funcionamento e empréstimos adiante.

\item {\bf ET's (Estações de Trabalho):} são salas com alguns terminais que têm
acesso à Internet e muito mais. Não, bixes, não se animem, pois ela não é para
vocês! Essa sala só pode ser usada pelos alunos da pós-graduação e por alguns
alunos que fazem iniciação científica.

\item {\bf Rede Linux:} podem ser usadas por alunos da graduação (ieiii!).
Utilizam o Linux, que, ao contrário do Windows, é um Sistema Operacional.

\end{itemize}

\end{subsecao}

\begin{subsecao}{Bloco B}


No bloco B, vocês devem conhecer as salas de aula (da graduação), a Vivência, as
mesas azuis (que não são todas azuis), o CEC, a seção de alunos, a lanchonete, 
o CAEM e o Xerox para vocês copiarem \sout{o caderno do colega} o capítulo do 
livro para estudar para a prova.

\begin{itemize}
\item {\bf Vivência:} sala 18, arrasadora de graduações, é onde as pessoas podem
dormir, assistir TV, jogar sinuca, cartas, xadrez, vídeo-game e até
mesmo estudar. Guardem bem esse nome, vocês vão passar a maior parte do tempo lá.
Lá também estão os armários e as salas do CAMat e da Atlética.

Os armários são renovados semestralmente. Após o período de renovação os que
estiverem sobrando ficam livres para serem alugados. Aí cabe a vocês, bixes, ficarem
atentos às datas divulgadas pelo CAMat e demonstrarem interesse pelos armários para
conseguirem entrar no sorteio.

\item {\bf Lanchonete:} uma alternativa para os dias chuvosos que você está sem
  guarda-chuva para ir até o bandejão é a lanchonete do IME. Tem um vasto
  cardápio de salgados e congelados esquentados no microondas, além de sucos,
  refrigerantes, chocolates, bolachas e similares.

\item {\bf CEC (Centro de Ensino de Computação):} é um centro munido de dezenas
de computadores rodando Windows, e alguns rodando Linux. Também é possível fazer
(ou pelo menos tentar) EPs em situações de desespero (vocês vão descobrir o que é
isso logo, logo). Algumas aulas de computação são ministradas lá (como Desafios
de Programação do BCC).

\item {\bf Seção de Alunos:} é aqui que vocês vão resolver quase todos os
pepinos de vocês durante a Graduação. Desde descobrir por que o Júpiterweb não
quer matricular vocês em uma disciplina obrigatória até fazer a
MATRÍCULA PRESENCIAL nos dias 27/02 e 28/02. %REFTIME

A Seção de Alunos fica na sala 12, ao lado do CEC e das mesas azuis.
\begin{itemize}
\item[-] Horário de atendimento: 10:00 às 13:00 e 17:00 às 20:00, de segunda a sexta.
\item[-] E-mail: \tt{saol@ime.usp.br}
\item[-] Telefone: \tt{11 3091-6149} e \tt{11 3091-6279}
\end{itemize}

\item {\bf CAEM:} sigla para o Centro de Aperfeiçoamento do Ensino da
  Matemática. É um órgão de extensão, que oferece cursos, oficinas, palestras e
  presta serviços de assessoria para professores de Matemática. Vocês que fazem
  Licenciatura (futuros professores) também podem usufruir dos serviços do CAEM,
  e até mesmo estagiarem lá. Fica no 1º andar, em frente às escadas.

\item{\bf Xerox:} fica no térreo, saindo da vivência e indo até o final do segundo
  corredor à direita (o que não tem os banheiros no final). Para fazer cópias,
  é preciso primeiro comprar fichas na salinha da Atlética (que fica na Vivência).
\end{itemize}

\end{subsecao}

\mapaime{bloco_A_terreo}

\mapaime{bloco_A_piso1}

\mapaime{bloco_A_piso2}

\mapaime{bloco_B_terreo}

\mapaime{bloco_B_superior}

\begin{subsecao}{Bloco C}

Apesar de atualmente abrigar os professores e a secretaria do departamento de
Computação, o misterioso Bloco C é um local a que nós, estudantes, infelizmente não
temos livre acesso. Para garantir tranquilidade e boas condições de trabalho aos
docentes, os alunos devem ser anunciados ou apresentarem uma boa desculpa para
adentrar o local. (Na verdade, os professores do MAC acham que os alunos são monstros
verdes e gosmentos, com $e^{10}$ braços, $\pi$ olhos e que comem criancinhas;
por isso, não querem esse tipo de criatura perambulando pelos seus corredores.)

\end{subsecao}

\begin{subsecao}{Bloco D (vulgo C')}

O bloco D não passa de uma extensão do bloco C. Na verdade, eles são o mesmo bloco,
só que têm entradas independentes. Pertence ao NUMEC, que ninguém sabe ao certo
o que significa. Alguns dizem que ele não existe, e que é apenas fruto da sua
imaginação.

\end{subsecao}

\begin{subsecao}{CCSL}

O anexo do Bloco C abriga o CCSL --- Centro de Competência em \textit{Software}
Livre. O que tem lá? Algumas salas de docentes (afinal, ele é uma extensão do
Bloco C), alguns laboratórios (de sistemas, de inteligência artificial,
de visão computacional, de computação musical...), um pequeno auditório
onde rolam palestras e seminários e, finalmente, o laboratório de extensão,
que é usado para cursos rápidos e palestras com parte prática onde o público
precise de um computador. É esse laboratório que abriga grande parte dos grupos
de extensão do IME que vocês verão na seção ``Atitude, bixes!''.

Se vocês se interessam por \textit{Software} Livre, não deixem de ler a seção
do CCSL desse Guia!

\end{subsecao}

\begin{subsecao}{Biblioteca e SIBi-USP}

A Biblioteca Carlos Benjamin de Lyra é uma das 48 bibliotecas que compõem o Sistema Integrado de Bibliotecas (SIBi). Considerada uma das maiores bibliotecas na Área de Matemática da América Latina, seu acervo é formado por livros, relatórios, revistas acadêmicas, teses, dissertações, trabalhos de docentes, CDs e DVDs totalizando cerca de 227.000 itens. 

\begin{subsubsecao}{Instalações:}

O espaço físico da biblioteca é de 1.642m\textsuperscript{2}. Possui um salão de estudos (mesinhas individuais) com funcionamento 24 horas, 13 salas para estudo em grupo (pra conseguir usá-las é preciso pegar a chave no balcão de empréstimos) e bancadas para estudo individual, totalizando 220 assentos.
\end{subsubsecao}

\begin{subsubsecao}{Empréstimos:}

O aluno de graduação poderá emprestar 10 livros por 10 dias. Para o empréstimo é necessário cadastrar uma senha pessoalmente e apresentar o cartão provisório, ou carteirinha USP, ou aplicativo e-card, ou RG/CNH. Poderão ser feitas 3 renovações pela internet (DEDALUS) e 3 reservas de livros que estiverem emprestados com outros alunos. O atraso na devolução gera suspensão (nº livros atrasados multiplicado pelo nº dias em atraso). 
\end{subsubsecao}

\begin{subsubsecao}{Recursos para pesquisas:}
\begin{itemize}
    \item DEDALUS - é o catálogo online para consulta e localização dos materiais nos acervos das bibliotecas da USP e para renovação de empréstimos e pedidos de reservas. 
    \item Portal de Busca Integrada - interface única para pesquisa de documentos impressos das bibliotecas e dos documentos disponíveis online com textos completos. 
    \item Bases de Dados e Portal Capes – acesso online a e-books e artigos de revistas assinadas pela USP e pela Capes.
    \end{itemize}
\end{subsubsecao}

\begin{subsubsecao}{Serviços oferecidos:}

Além dos empréstimos, a Biblioteca possui scanner de autoatendimento e oferece os seguintes serviços para os alunos:
\begin{itemize}
    \item COMUT (comutação bibliográfica): obtenção de cópias de artigos e capítulos de livros nas principais bibliotecas brasileiras; 
    \item EEB (empréstimo entre bibliotecas): empréstimo entre bibliotecas de instituições externas à USP ou bibliotecas USP interior;
    \item Orientação para trabalhos acadêmicos: orientação para formatação de TCC, teses, dissertações, citações e referências bibliográficas;
    \item Treinamentos e workshops: uso dos recursos da biblioteca e pesquisas em portais e bases de dados;
    \item Apoio ao pesquisador: auxílio na localização de artigos, livros e referências bibliográficas;
    \item Aplicativo Bibliotecas USP: pesquisa no acervo das bibliotecas da USP diretamente do Android, iPhone, iPod Touch ou iPad; 
    \item Conexão VPN: acesso remoto (de casa) às bases de dados, revistas eletrônicas e e-books.
\end{itemize}
\end{subsubsecao}

\begin{subsubsecao}{Horário de funcionamento:}

Período letivo: de 2ª a 6ª feira das 8h00 às 21h30 \newline
\phantom{Período letivo: }de sábados das 9h00 às 13h00

Período não letivo: de 2ª a 6ª feira das 8h00 às 20h00 \newline
\phantom{Período não letivo: }de sábados fechada 
\end{subsubsecao}
\begin{description}
  \item[Facebook:] \url{https://www.facebook.com/bibliotecaimeusp/}
  \item[Site:] \url{https://www.ime.usp.br/bib}
  \item[E-mail:] \texttt{bib@ime.usp.br}
  \item[Telefone:] 3091-6109 e 3091-6174
\end{description}

\end{subsecao}
\end{secao}


\quadrinhos{7}

% Atitude, bixo! ---------------------------------------------------------------
\begin{secao}{Atitude, bixes!}

Na USP, os alunos têm a liberdade e apoio de se organizarem
para montar grupos de debates, ciclos de palestras, grupos
de desenvolvimento e até mesmo grupos para jogarem alguma coisa (como
RPG, Magic, Yu-Gi-Oh ou algum esporte).
Portanto, caso vocês tenham algum projeto em mente, não hesitem
em se organizarem com seus amigos e se informarem em como avançar com essa
ideia. Lembre-se de que seus veteranes estão aí para te aconselhar e tirar
eventuais dúvidas.

É possível também juntar-se com alguns amigos e formar grupos de
estudo, seja para alguma matéria com a qual vocês tenham dificuldade, para
discutir aquele EP/lista de exercícios que ninguém está conseguindo
fazer ou simplesmente estudar algum tópico de interesse mútuo.

A seguir, temos (majoritariamente) alguns dos exemplos dos grupos que foram 
direta ou indiretamente criados por alunos do nosso Instituto!

% RDs --------------------------------------------------------------------------
\begin{subsecao}{RDs}

Antes de mais nada, RD significa Representante Discente. O RD é um aluno que
representa nossos interesses frente aos diversos conselhos e comissões
existentes, sendo um forte elo de ligação entre professores e alunos. O RD
ajuda a tomar decisões que impactam todo o IME, como autorização para festas,
mudanças no currículo, aumento de vagas na FUVEST, quantidade de bolsas,
reformas, mudança no corpo docente (às vezes lutamos para tirar algum
professor), enfim, coisas desse tipo e muitas mais.

Acho que já deu para perceber o quanto é importante ter um estudante em cada um
desses conselhos. Infelizmente, não costumamos preencher todas as vagas. Isso se
deve ao desinteresse de alguns ou falta de tempo da maioria de seus veteranes.
E vocês são quem tem mais tempo para fazer as coisas funcionarem aqui, já que
ainda não sabem o que é rec, DP, trabalho, estágio etc. Portanto, se quiserem
fazer alguma coisa pelo lugar onde vocês, bixes, vão estudar, está aí
uma dica.

Como RDs vocês poderão entender melhor o funcionamento do Instituto e ajudar no
processo de melhorá-lo. Vocês também poderão melhorar o relacionamento entre
estudantes, professores e servidores e entender como os professores pensam.
As eleições são organizadas pela direção do instituto no final do segundo
semestre e o mandato é de um ano.

Aqui vai um breve resumo do que mais ou menos acontece em cada um dos colegiados
nos quais temos direito a representante(s):

No IME, temos 24 cargos de RD com 34 vagas no total, sendo 17 reservadas
para graduação, 14 para pós e 3 livres. Todos têm direito a um suplente.
Caso o RD não possa ir a alguma reunião ou, por algum motivo da vida, tenha
que abandonar o cargo, o suplente assume em seu lugar e o cargo não fica
sem representante.

Existem diferentes níveis de hierarquia na administração.

{\bf As CoCs,
Comissões Coordenadoras de Curso (Lic, Pura, Estatística, Aplicada, BMAC e
Computação)} são as mais próximas dos alunos. Temos um cargo de aluno em cada
comissão. São comissões pequenas, que tratam dos problemas internos de cada
curso: mudança de currículo, requerimentos, optativas, etc. São subordinadas
à CG e ao conselho do relativo departamento. Analogamente, temos um cargo em cada
Comissão Coordenadora de Programa (de Pós).

{\bf Os Conselhos de Departamento (MAT, MAE, MAC e MAP)} têm uma dinâmica um
pouco diferente das CoCs: são mais formais. Cada conselho se reúne (quase)
mensalmente e são formados (em geral) por mais pessoas, sendo que existem
regras sobre participação dos diferentes níveis hierárquicos de
professores (Titular, Associado, Doutor e Assistente). Nesses conselhos, além
de aprovar algumas das decisões das Comissões Coordenadoras de Curso e de
Programa (pós) e distribuição de carga didática, são discutidos re-oferecimento
de curso, revisão de prova, supervisão das atividades dos docentes,
afastamentos (temporários ou não), contratação de professores e muitas outras
coisas.
Os Conselhos de Departamento são subordinados à Congregação e ao CTA.

{\bf A Comissão de Graduação (CG)}, basicamente, avalia requerimentos,
mudança/criação de cursos e jubilamentos. Analogamente, existe a Comissão de
Pós-Graduação (CPG). Ambas são subordinadas à Congregação.

{\bf A Comissão de Cultura e Extensão (CCEx)} quase nunca tem reunião. Cuida
das atividades de extensão: Matemateca, CAEM etc.

Também há comissões mais específicas, como a comissão de estágio, a comissão de
pesquisa (do doutorado) e o Centro de Competência em Software Livre (CCSL), da
computação.

Os dois conselhos mais importantes são o CTA e a Congregação, ambos presididos
pelo Diretor.

{\bf O Conselho Técnico e Administrativo (CTA)} cuida de todas as questões não
acadêmicas: Orçamento, reformas, avaliação dos funcionários, Xerox, etc. É
formado pelos quatro chefes de departamento, diretor, vice-diretor, um
representante dos funcionários e um RD.

{\bf A Congregação} é o órgão máximo do Instituto. Inclui muitos professores, a
maioria titular. São 3 RDs de graduação e 2 de Pós. Basicamente,
nesse órgão, são rediscutidas e aprovadas (ou não) muitas das decisões
dos órgãos subordinados. Os membros da Congregação têm voto na eleição para
Reitor e Vice-Reitor.

Bom, caso não tenha ficado claro desde o começo desse texto, percebam que é
muito importante ter um aluno em cada um desses conselhos. Se estiverem tendo
problemas com professores, requerimentos etc., ou simplesmente quiserem saber
o que anda acontecendo, procurem o RD certo pra conversar. Perguntem,
participem, votem e façam o IME um lugar melhor.

Sobre a eleição dos RDs: A eleição oficial para os RDs acontece no final do ano
(então fiquem atentos!) e é organizada pela diretoria do instituto. Os
interessados devem preencher um formulário de inscrição e levar até a
Assistência Acadêmica do IME, que vai organizar todos os inscritos e abrir um
processo de eleição online (em que todo IMEano pode votar). Quando a eleição
estiver se aproximando, vocês receberão (vários) e-mails com os documentos
necessários, prazos e links de votação.

%REFTIME
O resultado da eleição anterior com os RDs de 2019 pode ser encontrada no site
do IME:
\url{www.ime.usp.br/eleicoes-estatutarias}

\end{subsecao}


% Rede Linux -------------------------------------------------------------------
\begin{subsecao}{Rede Linux}

\figurapequenainline{rede_linux}

\begin{subsubsecao}{Introdução}

A Rede Linux é uma rede de computadores, administrada por alunos do IME e
que fornece diversos serviços para todos os alunos do IME.
Ela disponibiliza:

%FIXME
\vspace{-1em}

\begin{itemize}
\item 2 salas de computadores (no bloco A) com todo\footnote{ Se um programa
estiver faltando, mande um email pra admin@linux.ime.usp.br pedindo-o.} tipo de
programa necessário para suas atividades acadêmicas (com pelo menos uma que fica
aberta 24 horas por dia, 7 dias por semana\footnote{ Mas talvez vocês não
consigam entrar no bloco A depois da meia noite, que é quando a portaria
fecha.});
\item Uma página na internet para cada aluno;
\item Um \textit{e-mail} para cada aluno;
\item Espaço para você guardar seus arquivos;
\item Acesso remoto via ssh (linux.ime.usp.br);
\item Impressoras;
\item Admins dispostos e capazes, para o caso de algum usuário ter alguma boa
ideia para adicionar a esta lista;
\end{itemize}
\end{subsubsecao}

\begin{subsubsecao}{O Linux}

A rede utiliza em todos os seus computadores um sistema operacional chamado
Linux. Esse é um sistema desenvolvido de forma colaborativa pelos usuários
e empresas interessados nele (se quiser saber mais a respeito, pesquise
por ``software livre''!).

O Linux não é um sistema mais difícil de usar do que o Windows. É apenas
diferente em alguns aspectos. Além de tudo, existem cursos de Linux que são
organizados pelos alunos do IME. Os admins costumam promover esses cursos.
Fiquem atentos aos emails!

Não se deixem intimidar pelo sistema. Se vocês se derem ao trabalho de
aprender a utilizá-lo bem, verão que ele é bastante flexível, e até mesmo
interessante (tanto quanto um sistema operacional pode ser =P).

\end{subsubsecao}

\begin{subsubsecao}{Os admins}

Os admins são alunos do bacharelado em ciência da computação (vulgo BCC) que
são responsáveis por administrar a rede. Entre outras coisas, isso quer dizer
manter os computadores funcionando, ajudar os alunos a usar a rede (com
cursos\footnote{ Fiquem atentos aos emails!!!} e resolvendo dúvidas nos
horários de plantão\footnote{ Na página da rede, estão os horários de todos os
admins.}) e também implementar coisas novas na rede (aceitamos sugestões!)

Os admins são escolhidos por um treinamento que acontece de dois em dois anos,
em todo ano par. Mais informações serão divulgadas quando este estiver próximo
a ocorrer.

\end{subsubsecao}
\begin{subsubsecao}{Como criar uma conta?}

Basta passar na Admin, na sala 125 do bloco A (como vocês são bixes: bloco A é o da
biblioteca, bloco B aquele que tem muitas salas de aula e que vocês vão passar boa
parte da vida de vocês). Contatos:

%FIXME
\vspace{-1em}

\begin{description}
\item [e-mail:] admin@linux.ime.usp.br
\item [Página:] \url{www.linux.ime.usp.br}
\item [Sala:] 125, bloco A
\end{description}

%FIXME
\vspace{-.5em}

\end{subsubsecao}

\end{subsecao}


% FLUSP ------------------------------------------------------------------------
\begin{subsecao}{FLUSP}

\figurapequenainline{flusp}

Você já ouviu falar sobre Linux? E GCC? Talvez Gimp, qbittorrent, VLC ou Firefox?
Todos estes projetos são chamados de FLOSS: Free Libre Open Source Software, e
têm em comum a liberdade de código e conhecimento. Isto quer dizer que qualquer
um pode ler, usar e contribuir para estes projetos, seja com código ou arte!

O FLUSP: FLOSS@USP é um grupo de extensão com o objetivo de reunir alunos de
graduação e pós-graduação interessados em contribuir para projetos FLOSS.
Atualmente, temos contribuidores no Kernel Linux, no compilador GCC, no controlador 
de versão git, no projeto Caninos Loucos e muitos outros. Em 2019, o FLUSP 
foi responsável por aproximadamente 20\% das contribuições para o Kernel Linux no 
subsistema Industrial Input/Output. Devido a nossas contribuições para os drivers 
da Analog Devices, a empresa doou duas placas de testes ao grupo. Além delas, temos 
uma placa Labrador doada pelo projeto Caninos Loucos.

Assim como na comunidade FLOSS, nós encorajamos a liberdade no FLUSP. Seja para
contribuir com um projeto existente ou compartilhar um projeto pessoal para
outros contribuirem, aqui você encontra espaço!

\begin{description}
  \item[Facebook:] \url{facebook.com/flusp}
  \item[Telegram:] \url{http://tiny.cc/flusp}
  \item[IRC:] Servidor \texttt{irc.freenode.net}, canal \texttt{\#ccsl-usp}
  \item[Site:] \url{https://flusp.ime.usp.br}
  \item[Lista de email:] \texttt{flusp@googlegroups.com}
  \item[GitLab:] \url{https://gitlab.com/flusp}
\end{description}

\end{subsecao}


% IME Júnior -------------------------------------------------------------------
\input{ime_jr.tex}

% MaratonUSP -------------------------------------------------------------------
\begin{subsecao}{MaratonUSP}

\figurapequenainlineapertada{maratonusp}

O MaratonUSP é um grupo de estudos de programação competitiva. O principal evento é a Maratona de Programação, uma competição anual organizada pela Sociedade Brasileira de Computação (SBC). A competição é restrita a universitários que se organizam em times de três pessoas para resolver os mais variados problemas em forma de código. Apesar da coletividade, há uma dificuldade a mais: cada equipe dispõe de apenas um computador.

A Maratona é dividida em várias fases. Primeiro, há uma seleção interna dos times que representarão a USP na competição. Uma das vagas é reservada para um time de bixes! Depois, há uma fase regional com outras universidades da capital paulista, onde a USP tem um histórico de vitórias absoluto.

As coisas são consideravelmente mais difíceis na fase nacional, conhecida como Final Brasileira. As melhores universidades do Brasil se reúnem em algum canto do país para disputar as vagas brasileiras para a fase mundial. Exemplos recentes são Belo Horizonte, Foz do Iguaçu e Salvador. A hospedagem, comida e transporte são todos cobertos pela organização e pela universidade. Viagem de graça! Os melhores times brasileiros vão para a Final Mundial, sempre num país diferente. Nós estivemos recentemente no Marrocos, na Tailândia e na China. Em hotéis 5 estrelas! De graça!

Grandes empresas têm muito interesse nos alunos que participam da Maratona. É muito comum encontrar ex-maratonistas em empresas como Google, Facebook, etc. O pensamento abstrato e o trabalho em equipe desenvolvidos na Maratona são características muito bem avaliadas pelos recrutadores.

O ano começa com o bixeCamp, uma série de aulas focadas em ensinar o básico de programação e algoritmos. As aulas são no início dos treinos, que ocorrem no CEC (do lado da Seção de Alunos) à partir das 14h às sextas-feiras. O treino acaba às 19h, mas não é necessário ficar até o final. Alunos já experientes em programação são encorajados a comparecer também. Te esperamos lá!

\begin{description}
\item [Facebook:] \url{facebook.com/MaratonUSP}
\item[Site:] \url{www.ime.usp.br/~maratona}
\end{description}

\end{subsecao}


% USPCodeLab -------------------------------------------------------------------
\input{usp_code_lab.tex}

% USPGameDev: Pesquisa e Desenvolvimentos de Jogos na USP ----------------------
\begin{subsecao}{USPGameDev: Pesquisa e Desenvolvimentos de Jogos na USP}

\figurapequenainline{uspgamedev}

Valve. Blizzard. Rockstar. Nintendo. USPGameDev. O que esses nomes têm em comum?
São nomes de grupos de desenvolvedores de jogos. E um deles tem sua sede na USP.

Constituído primariamente, mas não exclusivamente, de alunos da USP de diversas
áreas (computação, matemática, design, música, letras, \textit{etc}.) o
USPGameDev (UGD) foi criado em 2009  e já lançou dezenas de jogos (um deles no 
Steam!\footnote{\texttt{https://store.steampowered.com/app/827940/Marvellous\_Inc/}}) 
e até mesmo o seu próprio \textit{kit} de desenvolvimento. Trabalhamos com jogos 
digitais e analógicos (video-games e jogos de tabuleiro, por exemplo). Vale 
ressaltar que adotamos a filosofia de software livre (\textit{``livre'' de 
``liberdade'', não necessariamente grátis}).

Queremos aprender e ensinar desenvolvimento de jogos como a atividade 
multifacetada que ela é. Você, que acabou de ingressar, também pode criar o seu 
próprio jogo dentro do USPGameDev, além de participar dos muitos projetos que já 
estão em andamento. Tanto que alguns dos nossos últimos jogos foram produtos de 
grupos de ingressantes (com muito tempo livre ou em disciplinas que aceitavam 
jogos como trabalhos). 

Não é necessário conhecimento prévio algum! Todos nós começamos a vida sem saber 
programar, desenvolver, projetar, desenhar, \textit{etc}. Isso porque não somos 
uma \textit{empresa} de jogos, mas um grupo de estudos. E, portanto, buscamos 
aprender o que não sabemos e ensinar o que sabemos.
Justamente por isso, o UGD também oferece cursos e \textit{workshops} livremente 
para a comunidade USP sobre diversos assuntos envolvendo desenvolvimento de 
jogos. Fiquem de olho! Além disso, participamos de \textbf{game jams} (ou 
hackathons): eventos regionais e internacionais onde temos de 24 a 72 horas para 
fazer um jogo com base em um tema que só é revelado na hora!

Interessados? Acessem nosso muito bem desenvolvido \textit{site} e deem uma 
conferida na página de downloads! 

\textbf{Para participar, basta entrar em contato (via email ou o que seja) 
conosco e marcar um dia para conversarmos}. Ser um membro não é lá muito formal, 
a gente bate um papo e vê o serial legal para você fazer. O grupo é horizontal e 
cada um escolhe quanto participa.

\begin{description}
  \item[Site:] \url{http://uspgamedev.org}
  \item[Fórum:] \url{forum.uspgamedev.org}
  \item[E-mail:] contato@uspgamedev.org
  \item[Facebook:] \url{facebook.com/UspGameDev}
\end{description}

\end{subsecao}


% IMEsec -----------------------------------------------------------------------
\begin{subsecao}{IMEsec}

\figurapequenainline{imesec}

O IMEsec é um grupo de extensão focado em aprender, estudar e se divertir com a
segurança da informação. Sem que a maioria das pessoas se dê conta, este nicho
está presente no cotidiano. Mandar uma mensagem no WhatsApp, navegar pela web,
entrar no Facebook: exemplos de ações simples do dia-a-dia que necessitam ser
feitas de maneira segura, visando à privacidade do usuário.

Nosso grupo, formado majoritariamente por alunos do IME-USP no início de 2017,
busca entender melhor este universo e expandí-lo no ambiente universitário. O
foco tem sido amplo; desde resolução de desafios on-line (que são muito
divertidos), participação em competições (sim, competições — \textit{capture the
flag} — muito, muito legais), apresentação de palestras e até desenvolvimento de
projetos que possam beneficiar a população. Todos são bem-vindos; basta ter
interesse pelo assunto.

As experiências decorrentes de nossas atividades ajudaram e ajudam nos estudos,
além de agregarem valor à graduação. Por meio desses, mergulhamos não só em
computação mas também em matemática, estatística e até outras áreas bem
inusitadas.

Fiquem de olho na nossa página do Facebook e do Telegram para saber mais sobre
as reuniões semanais, além de fatos interessantes (ou simplesmente engraçados)
sobre o vasto mundo da segurança. Participem!

\begin{description}
  \item[Facebook:] \url{facebook.com/imesec}
  \item[Telegram:] \url{tiny.cc/imesec-telegram}
  \item[Site:] \url{imesec.ime.usp.br/}
\end{description}

\end{subsecao}


% Hardware Livre ---------------------------------------------------------------
\input{hardware_livre.tex}

% Tecs: Computação Social ------------------------------------------------------

\begin{subsecao}{Tecs}

\figurapequenainline{tecs}

O Tecs é um grupo de extensão que iniciou suas atividades no segundo semestre de
2017, e é focado no impacto social da computação e da tecnologia. 

Buscamos construir uma rede de estudantes interessados em melhorar a comunidade
local e a sociedade e criar oportunidades para que eles possam engajar-se em
projetos com esse fim; unir esforços para formar uma sociedade e profissionais
éticos e conscientes sobre o uso da tecnologia; promover a educação tecnológica
igualitária da população por meio de cursos, oficinas e ações promovidas pelo
grupo; e estimular alunos a usarem a tecnologia para solucionar problemas da
comunidade. 

Em termos gerais, pretendemos que os estudantes entendam como a tecnologia pode
ser utilizada para o bem coletivo, e utilizem esse conhecimento na prática, por
meio de colaborações com a comunidade local, os serviços públicos, as
organizações não-governamentais e as sem fins lucrativos. 

Se você tem interesse em promover o ensino de computação, em debater questões
éticas e sociais no contexto tecnológico ou em desenvolver aplicativos, sites ou
sistemas em parcerias com projetos sociais, entre em contato conosco e participe
do grupo!

\vspace{-1em}
\begin{description}
  \item[Site:] \url{www.ime.usp.br/~tecs}
  \item[Facebook:] \url{www.facebook.com/tecs.usp}
  \item[Telegram:] \url{https://t.me/tecsusp}
\end{description}

\end{subsecao}


%FIXME GAMBIARRA para não quebrar página num lugar zuado
\pagebreak

% Diversime --------------------------------------------------------------------
\input{diversime.tex}

% Existimos! -------------------------------------------------------------------
\begin{subsecao}{$\exists$xistimos!}

\figurapequenainlineflexivel{existimos}{30pt}


$\exists$xistimos! surgiu em 2014 com a proposta de criar um espaço de
confiança entre as alunas do IME, onde cada uma possa ter a liberdade e
segurança para discutir suas vivências e propostas.


O desconforto com a forma como as mulheres são tratadas no IME e nas ciências
exatas em geral fez com que nos juntássemos para discutir como as questões de
gênero se manifestam no instituto e quais são as formas de agir para evitar
situações desconfortáveis ou preconceituosas.

Desde então promovemos uma série de eventos e intervenções para que tal debate
atinja toda a comunidade imeana, além de realizarmos reuniões periódicas apenas
com mulheres para que possamos conversar e nos ajudar, em qualquer situação,
mas em especial naquelas onde possamos ser vítimas ou testemunhas de  
preconceito e machismo. Todas vocês estão convidadas a participar das nossas  
reuniões, alunas de outros institutos são muito bem vindas $<$3.

Para falar conosco ou participar do grupo basta enviar um email ou entrar em
contato pelo facebook:


\begin{description}

\item[E-mail:] 3xistimos@gmail.com ou existimos@google.groups.com
\item[Facebook:] \url{https://www.fb.com/3xistimos/}
\item[Para denúncias ou relatos anônimos acesse:] \url{http://bit.ly/existimos}

\end{description}


\end{subsecao}



%FIXME GAMBIARRA para não quebrar página num lugar zuado
\pagebreak

% Comissão de Acolhimento da Mulher! -------------------------------------------
\input{cam.tex}

%FIXME GAMBIARRA para não quebrar página num lugar zuado
\pagebreak

% CinIME -----------------------------------------------------------------------
\begin{subsecao}{CinIME}

\figurapequenainline{cinime}

Convidamos vocês a conhecer o CinIME!! Somos um projeto do CAMAT que promove o
lazer e a integração de IMEanos, entre docentes, discentes e funcionários por
meio de sessões de cinema em que vocês escolhem o filme!

O CinIME ocorre toda sexta-feira, às 16h, na sala B5. A sessão, o refrigerante
e a pipoca são grátis. Funciona assim: você entra na sala, se serve do que
quiser, senta, relaxa e se diverte assistindo ao filme da semana ao lado dos
seus amigos.

A escolha do filme é feita através de votação no grupo “CinIME USP”, no
facebook. Para participar, primeiro você acessa e curte a página do CinIME no
facebook (https://www.facebook.com/cinimeusp) para ter acesso às nossas
atualizações. Então, na aba “Sobre”, você clica no link de sugestões apontado
na descrição da página. Dentre todas as sugestões de filmes, cinco são
escolhidos pela equipe do CinIME para a votação, priorizando a diversidade de
gêneros cinematográficos. Depois disso, a votação é lançada no grupo “CinIME
USP”, e o resultado é sempre anunciado ao fim da sessão da semana em curso, que ocorre
na sexta à tarde. O filme eleito é divulgado nos grupos
dos bixes e no grupo “IME-USP”, no facebook, e também por cartazes espalhados
pelo instituto. Uma vez por mês, a equipe do CinIME toma a liberdade de
escolher o filme da sessão sem votação.

Não se esqueçam de sugerir filmes, votar e comparecer ao CinIME!

Curtam nossa página no facebook e fiquem por dentro das novidades!

Parabéns e boa sorte nesta nova jornada, bixes!

\begin{description}
  \item[Facebook:] \url{https://www.facebook.com/CinIMEUSP}
\end{description}

\end{subsecao}


% Grupo A5 ---------------------------------------------------------------------
\begin{subsecao}{Grupo A5}

\figurapequenainline{grupo_A5}

Somos um grupo de estudantes da graduação e pós-graduação do IME e, anualmente,
organizamos eventos acadêmicos no instituto. Tudo começou em 2012, quando dois
alunos da pura perceberam que alguns assuntos muito interessantes sobre
matemática e sobre o meio acadêmico nem sempre eram desenvolvidos em sala de
aula, e resolveram organizar um evento para levar um destes assuntos aos demais
estudantes instituto. Assim, juntamente ao CAMat, organizaram um ciclo de
palestras sobre "Os 7 Problemas do Milênio", e o evento foi um sucesso.  Vários
professores e estudantes começaram a pedir que mais eventos como este fossem
organizados, e foi então que um deles teve a ideia de criar um grupo
independente das demais instituições do IME (sim, o Grupo A5 é independente. Não
é vinculado ao CAMat e nem a nenhum outro grupo, apesar de aceitar parcerias em
alguns eventos), com a finalidade de complementar a formação dos estudantes do
IME e de quem quiser participar, levando palestras e outros eventos sobre temas
relevantes e não explorados no currículo, de forma gratuita e com linguagem de
fácil entendimento.  E assim, no final de 2013, o Grupo A5 oficialmente nasceu,
com nome e logo e com novos integrantes no grupo.  Desde então, não paramos
mais. Em 2014 organizamos o ciclo de palestras "IC ou Não IC? - Eis a Questão",
em 2015 organizamos o evento "História da Matemática", que foi indicado como um
dos destaques de Cultura e Extensão de 2015 pelo IME. E em 2018 realizamos, em 
parceria com o Existimos, o ciclo de palestras "Mulheres no Mundo Corporativo",
que foi um sucesso! 

Para mais informações do Grupo A5 e dos eventos já organizados por nós, acessem
nosso site \url{www.ime.usp.br/~acinco} (ainda está em construção, mas em breve os
vídeos das palestras e demais informações estarão lá). E não deixem de curtir
a página Grupo A5 no facebook: \url{www.fb.com/pagina.GrupoA5} (é aqui que
vocês terão em primeira mão os detalhes de tudo o que for feito por nós).

Vale ressaltar que o Grupo A5 é formado e mantido por estudantes do IME, então o
sucesso e a continuidade do grupo dependem de todos; começando por vocês,
bixes. Então venham, assistam, dêem sugestões, participem. E se gostarem,
juntem-se ao nosso grupo!

\end{subsecao}


% Olimpíadas de Matemática e Informática ---------------------------------------
% Maratona de Programação ------------------------------------------------------
\begin{subsecao}{Olimpíadas de Conhecimento}

\begin{itemize}

\item{\bf Matemática: }

Bom pessoal, se vocês entraram no IME, muito provavelmente já participaram
de alguma Olimpíada de Matemática no Ensino Fundamental e/ou Médio. A
boa notícia é que vocês vão poder continuar participando se quiserem,
e quem nunca participou tem a oportunidade de começar agora.

Mas por que participar? As Olimpíadas Universitárias de Matemática são uma
oportunidade de se divertir resolvendo problemas difíceis de Matemática e agregar
valor ao currículo ao mesmo tempo. Elas são parecidas com as Olimpíadas de
Ensino Médio, mas com conteúdo de Matemática da graduação (essencialmente
Cálculo, Análise, Álgebra Linear, Álgebra, Combinatória e Teoria dos Números),
mas com enfoque em problemas que exigem criatividade e técnicas mais inovadoras,
muitas das quais vocês provavelmente não verão durante toda a graduação.

De quais olimpíadas podemos participar? Como alunos de graduação, vocês podem
participar da Olimpíada Iberoamericana de Matemática Universitária (OIMU),
Olimpíada Brasileira de Matemática (OBM) e Olimpíada Internacional de
Matemática (IMC).

Como fazemos para nos preparar? Os sites institucionais dessas olimpíadas
têm todo o material necessário para vocês que querem estudar e se preparar
para elas.

Como fazemos para participar? Inscrevam-se pelo site ou entrem em contato com
o professor Yoshiharu. Para o IMC aconselha-se ter ganhado medalha na OBM,
já que é necessário apoio financeiro do IME por ser uma olimpíada internacional.

%REFTIME
Mas nós, bixes, temos chance? Como foi o desempenho de IMEanos nelas? Nós
obtivemos sucesso nestas olimpíadas. Ganhamos medalhas em todas as três
competições e o resultado mais recente foi uma medalha de bronze na IMC e ouro
na OBM.

Se tiverem alguma dúvida, não hesitem em perguntar a algum veterane sobre os
Olímpicos!

Links institucionais:

\begin{description}
  \item[] \url{http://www.cimat.mx/Eventos/oimu/}
  \item[] \url{http://www.imc-math.org/}
  \item[] \url{http://www.obm.org.br/opencms/}
\end{description}

\item{\bf Informática: }

\textit{``Informática? Vocês mexem com Word, Excel e PowerPoint então?''}

Responder essa pergunta já virou rotina para competidores da Olimpíada
Brasileira de Informática (OBI). Não, Informática não é Word. Oras, então o que é a OBI?

A OBI é uma competição de lógica, matemática e computação. As provas envolvem
alguns problemas que você deve resolver com programas de computador.

Esta competição, na graduação, é exclusiva para ingressantes recém formados do
ensino médio. Quer dizer que vocês são a nossa única esperança de trazer mais
gloriosas medalhas ao IME! Isso também quer dizer que essa é a sua única chance
de participar da OBI, uma competição relativamente tranquila comparada à
Maratona de Programação.

Para participar, basta falar com o MaratonIME
(\url{https://www.ime.usp.br/~maratona/}), um grupo de extensão focado nesse
tipo de competição que promete te ajudar a se inscrever e se preparar, ou com o
Professor Carlinhos (\url{http://www.ime.usp.br/~cef/}).

Para mais informações, acessem \url{http://olimpiada.ic.unicamp.br/}.

\item{\bf Maratona de Programação: }

À primeira vista, a Maratona de Programação pode soar um tanto
surreal. Nerds correndo pela USP ao mesmo tempo que resolvem
problemas de computação e matemática? Infelizmente esse não
é o caso.

A Maratona de Programação se resume à resolução de problemas.
Se você adora resolver desafios, quebrar a cabeça com novos
e excitantes problemas e acumular toneladas de dinheiro, esse
é o lugar perfeito para você!

A competição consiste em uma série de problemas que englobam
temas como programação dinâmica, grafos e estruturas de dados.
Times de três pessoas devem resolver a maior quantidade de
desafios em cinco horas de programação. E tudo isso com direito
a um lanche gratuito durante a prova.

Mas não temam, bixes. Não é só por que vocês acabaram de entrar que
a probabilidade de ganhar uma medalha seja nula. Inclusive, na primeira
fase da maratona, uma equipe de bixes tem vaga garantida para a
fase brasileira.

Além da fama, constantes pedidos por autógrafos e dinheiro de sobra,
a maratona também vai lhes trazer um conhecimento muito mais
adiantado em relação ao dos seus colegas de classe, e até oportunidades
de emprego em empresas de renome, como Google, Facebook e IBM.

Se vocês se interessaram pela maratona e querem saber os horários dos
treinos, como participar ou saber mais, acessem:

\begin{description}
  \item[Site:] \url{http://www.ime.usp.br/~maratona}
  \item[Site da competição nacional:] \url{http://maratona.ime.usp.br/}
\end{description}

\end{itemize}


\end{subsecao}


%FIXME GAMBIARRA para não quebrar página num lugar zuado
\pagebreak

% Fala Sério -------------------------------------------------------------------
\begin{subsecao}{Fala Sério}

\figurapequenainline{fala-serio}

A iniciativa Fala Sério é um grupo de extensão de cunho social da USP que atua em
escolas públicas de São Paulo proporcionando reflexões sobre futuro e felicidade.
O projeto se iniciou na Escola Politécnica da USP e devido ao seu grande sucesso 
se tornou um projeto de extensão da USP.

Para muitos alunos o processo de transição do Ensino Médio para a vida adulta é 
muito difícil e muitas vezes perturbador. Surgem muitas dúvidas, insegurança, 
cobranças; e é neste contexto que o Fala Sério atua.

Por meio de dois encontros levamos reflexões e espaços de conversa sobre 
felicidade e futuro; quais nossos sonhos e como alcançá-los. Desse modo, 
oferecemos um apoio pessoal que muitas vezes os jovens não encontram em suas 
escolas e que faz uma grande diferença em suas vidas.

Eaí, se animou com a ideia? Quer compartilhar suas experiências e entrar nessa 
missão com a gente? Visite nossa página no facebook e vem nos conhecer:
\url{facebook.com/iniciativafalaserio}

\end{subsecao}


\end{secao}


% CCSL - Centro de Competência em Software Livre
\begin{secao}{CCSL - Centro de Competência em Software Livre}

Se vocês já se cansaram da definição tradicional de software (``hardware é a
parte que você chuta, software é a parte que você xinga'') e prefere pensar que
``software é a parte com que você faz o que quiser'', o Software Livre é para
vocês. Baixar de graça? Sim! Copiar para os amigos? Claro! Fuçar, mexer e
hackear? Com certeza! E aqui no IME ele faz parte do ensino e da pesquisa,
literalmente, há décadas! Tanto que o IME tem o CCSL — Centro de Competência em
Software Livre — que existe para dar apoio a essas atividades internamente e
promover a difusão do Software Livre fora da Universidade também.

Vocês podem participar das atividades (palestras, eventos) que o CCSL
promove e também colaborar com elas; ou podem mergulhar em algum dos
projetos apoiados por ele, melhorando o código. Alguns desses projetos
podem vir a ser ótimos temas para seus TCCs ou para trabalhos de
iniciação científica (com bolsa! :D). Então vocês ainda juntam o útil ao
agradável. Visitem http://ccsl.ime.usp.br e inscrevam-se na lista de discussão,
conheçam os projetos e fiquem por dentro das próximas atividades do
Centro!

Ah, mas vocês ainda nem sabem direito o que é Software Livre? Não se preocupem.
Vocês ainda vão ouvir falar muito dele, tanto no IME quanto fora dele. %Além
%disso vai haver uma apresentação a respeito no começo do semestre, fiquem
%ligados! %REFTIME

\end{secao}


% CRInt ------------------------------------------------------------------------
\input{crint.tex}

% Cursos de Verão --------------------------------------------------------------
\input{cursos_verao.tex}

% CEC --------------------------------------------------------------------------
\begin{secao}{CEC - Centro de Ensino de Computação}

\begin{subsecao}{O que é?}

O Centro de Ensino de Computação (CEC) foi criado em 1993, pelo
Departamento de Ciência de Computação do IME, com o objetivo de
oferecer apoio às aulas dos cursos de graduação do Instituto e para os
treinamentos/cursos oferecidos à comunidade USP e não USP, através de
seus laboratórios com recursos de informática.
Está localizado no Bloco B do IME, ao lado da Seção de Alunos. O horário
de funcionamento do CEC é de 2a a 6a das 8h00 às 22h00 e nos recessos
escolares o horário de abertura é alterado para 9h00.

\end{subsecao}

\begin{subsecao}{Estrutura}

O CEC possui 137 computadores, distribuídos em 04 (quatro)
salas/laboratório, sendo 02 (duas) delas com capacidade para até 60
pessoas e com prioridade de uso para as aulas de graduação do IME e
para os cursos/treinamentos.
Nos outros laboratórios, os computadores são dual boot (linux e windows) e
possuem os aplicativos solicitados pelos docentes do Instituto. São
utilizados pelos alunos de graduação dos cursos do IME para a realização
dos trabalhos e pesquisas acadêmicos.

\end{subsecao}

\begin{subsecao}{Como funciona?}

Possui rede cabeada e para utilizá-la é necessário ter uma senha de
acesso (conta de usuário). Você precisa ser aluno matriculado no IME para
solicitar a senha de uso, pois o CEC não é uma sala pró-aluno (para ver
como as salas pró-aluno funcionam acesse www.usp.br/proaluno/). 
Para solicitar a senha de acesso, envie e-mail para cec-senha@ime.usp.br,
assunto: CEC senha, informe seu nome completo e número USP. Utilize
seu e-mail institucional (@usp.br - se esqueceu entre em
www.pedidoemail.usp.br).

Informações, como cota de disco/impressão, regras de utilização e faqs,
podem ser visualizadas em http://www.cec.ime.usp.br.

\end{subsecao}

Obs.: a partir de 2011, alguns pedidos de login são recebidos com maior simpatia
e resultam em logins e senhas iguais para Windows e Linux, ou seja, uma senha a
menos para você ter de lembrar.

Obs.2: bixes, ao frequentarem o CEC, fiquem atentos ao ar-condicionado. Se
estiver ligado, deem preferência a usarem calças, blusas, jaquetas e meias de
lã. Cobertores são opcionais. Se não, boa sorte ou \textit{hasta la vista}!

\end{secao}


% CAEM -------------------------------------------------------------------------
\begin{secao}{CAEM - Centro de Aperfeiçoamento do Ensino da Matemática}

O Centro de Aperfeiçoamento do Ensino da Matemática tem como principal
objetivo prestar assessoria aos professores de Matemática das redes
públicas (escolas municipais e estaduais) e particulares. Dentre
outras atividades, o CAEM oferece cursos, oficinas, palestras e
seminários voltados para professores de matemática dos níveis
Infantil, Fundamental e Médio. IMEanos podem usufruir e participar,
inclusive bixes. Aqui vão algumas informações úteis:

%Daqui para baixo é texto do próprio CAEM, com leves mudanças.

\begin{itemize}

\item \textbf{Como se cadastrar}: Para se cadastrar no CAEM a pessoa
  deve ter em mãos um comprovante de residência e 2 fotos
  3x4. Qualquer pessoa pode se cadastrar, independente se é aluno ou
  não.

\item \textbf{Materiais disponíveis}: O CAEM conta com um acervo
  voltado para o ensino da matemática nos níveis fundamental e médio,
  além de materiais de auxílio para aula como material dourado,
  sólidos, régua e compasso de lousa, DVD's. \textbf{Quais deles podem
    ser emprestados}: Livros e DVD's podem ser emprestados na sua
  maioria por um período de 1 semana (até 4 materiais, podendo
  renovar). Materiais de auxílio em sala podem ser emprestados em
  quantidade limitada e por apenas 2 dias. 
  
  Visite o CAEM, conheça os recursos e aproveite para conversar e
  tirar dúvidas com os colaboradores!   

\item \textbf{Informações sobre as oficinas}: Para todas as oficinas
  do CAEM são oferecidas apenas 5 vagas gratuitas a alunos do IME. Não
  são aceitos alunos além das 5 vagas, mesmo que esteja disposto a
  pagar, pois as oficinas e cursos do CAEM são voltadas a professores
  atuantes para trocas de experiência.

\item \textbf{Custos para o aluno}: Alunos do IME tem 5 vagas
  gratuitas nas oficinas.

\item \textbf{Horas de ATPAs}: Os certificados das oficinas e palestras
 do CAEM podem ser usados para as ATPA's.

\item \textbf{Horário de funcionamento}: Segundas e Quartas-feiras,
  das 10h00 às 19h00. Terças, Quintas e Sextas-feiras, das 10h00 às
  21h00.
\end{itemize}

Para mais informações, acesse o site \url{www.ime.usp.br/caem} ou veja
a página do CAEM no facebook procurando por ``CAEM IME USP''. Para
contato, envie um e-mail para {\tt caem@ime.usp.br}. Telefone e Fax:
(11) 3091-6160.


\end{secao}


% Um Pouco Sobre a USP ---------------------------------------------------------
\begin{secao}{Um Pouco Sobre a USP}

Vocês, que são novos na USP, devem saber desde cedo que aqui há muita
burocracia. É bom que se acostumem com ela, já que vocês terão que enfrentá-la.

O reitor é presidente do maior órgão da USP, o Conselho Universitário (abrevia-se
C.O. para evitar frases do tipo: ``Vou ter uma reunião no CU hoje'', ``O CU não
está funcionando muito bem esse semestre'', ``Os alunos não têm acesso ao CU''
etc.) que determina TODOS os rumos da universidade.

Abaixo dele vêm as coordenadorias e unidades. O SAS (Superintendência de Assistência Social),
por exemplo, é o departamento responsável pelos serviços
que a universidade oferece (não é a melhor coisa do mundo, mas oferece) para a
comunidade universitária: ônibus circulares, bandejões, moradia para estudantes
(CRUSP) etc.

Alguns outros lugares que vocês devem saber que existem são o HU (Hospital
Universitário), o CEPE (Centro de Práticas Esportivas, leia “cepê”), o banheiro
da FEA (Faculdade de Economia, Administração e Contabilidade), que fica na
frente do IME e é uma das faculdades mais bem abastecidas financeiramente
(apelidada “carinhosamente” de Shopping), a Física (certas aulas de laboratório são lá)
 e a querida Faculdade de Educação (para o pessoal da Licenciatura).

\begin{subsecao}{E-mail USP}

Bixes, atenção e cuidado ao e-mail que vocês receberam (ou vão receber)! É um
email@usp.br, do sistema “oficial” de e-mails da USP, o que significa que é
através dele que a Universidade, o Jupiterweb, a diretoria do IME, o CAMat e
alguns desocupados lhes enviarão comunicados oficiais, o que pode envolver desde
oportunidades diversas para intercâmbios, estágios, monitorias e tudo o mais que
vocês, simples bixes imeanos, sejam capazes de se imaginar fazendo. Por este e-mail
vocês receberão informações antes mesmo que as principais (é, nem todas são
colocadas em murais não) sejam também colocadas nos murais.

\end{subsecao}

\begin{subsecao}{Eduroam}

O Eduroam (EDUcation ROAMing, não tem nada a ver com alguém chamado Eduardo) é um
serviço de wi-fi que \sout{às vezes} provê acesso à Internet dentro de diversas
universidades (inclusive a USP). Mesmo que você seja aluno da USP, se viajar para
uma universidade no exterior que possui o Eduroam, ainda poderá ter acesso à Internet
usando o mesmo login!

Para usar o Eduroam, basta colocar o seu número USP no campo de login seguido por
@usp.br. Por exemplo, se o seu número USP for 123456789, então deverá colocar no
campo de login 123456789@usp.br. No campo da senha, deverá colocar a sua senha única
(a mesma senha que vocês usam para logar no Jupiterweb).

Tem várias outras informações sobre as configurações na hora de fazer login que
\sout{ninguém entende}, mas se algum de vocês entender, estão no site
https://eduroam.usp.br/como-usar/.

Para mais informações, acessem https://eduroam.usp.br/.

\end{subsecao}

\begin{subsecao}{O SAS}

A Superintendência de Assistência Social, que fica próxima à Praça do
Relógio, é o órgão da USP responsável pelo bem estar financeiro dos alunos — não
é lá a melhor coisa do mundo, mas ajuda... É onde vocês podem conseguir seus
milhares de benefícios, tais como auxílios financeiros, moradia gratuita (CRUSP)
ou bolsa alimentação (conhecido como vale-bandex) e dentistas gratuitos. Ou
seja, há bolsas de todo modelo, tamanho, design, estação e preferência
gastronômica nula que preferir. É também responsável pelo Setor de Passes
Escolares.

Se algum de vocês, bixes, estudou a vida inteira em escola pública, tem baixa renda,
sofreu agressões ou foi reprimido na infância, não pense que está sozinho nessa,
muitos também passaram por isso e vocês conseguem encontrar apoio aqui, permanência
estudantil é um direito e não deve ser motivo de vergonha, portanto busquem
auxílio,
pois podemos ajudar. Vão a seguir as diversas alternativas para todos que precisarem:

{\bf Moradia e auxílios financeiros}

Se vocês infelizmente não têm tanto acesso a meios de locomoção, ou dinheiro para
pagar transportes ou mesmo repúblicas, saibam que diversos auxílios podem ser
oferecidos para vocês para amenizar a sua situação:

Com a carteirinha provisória que você recebeu na matrícula, você deverá entrar
na página do SAS ({\tt http://sites.usp.br/sas/}) e solicitar a inscrição
para o processo de moradia e alojamento. Lá você terá um formulário, três
trilhões de coisas para preencher e assim que te chamarem, deve apresentar os
documentos (hum!) necessários para a assistente (você irá “ganhar” uma). Na
improvável hipótese do site estar fora do ar, como todo ano acontece, você terá
que ir ao SAS e pedir alojamento na USP, lá no Bloco G – sem trocadilhos - do
CRUSP.

A prioridade é dada aos residentes de outros estados ou interior de São Paulo,
mas os moradores de São Paulo também podem solicitar, dadas as proporções da
cidade e o tempo de duas horas para que moradores da Zona Norte venham às aulas.

Vale ressaltar, bixes, que vocês tomem o cuidado de não fazer o mesmo que outros
bixes burros de outros anos, que confundiram alojamento com moradia. São dois
requerimentos distintos e vocês deverão solicitar os dois se realmente quiserem
garantir um lugar para tomar banho e dormir.

Os requerimentos conterão perguntas sobre a situação socioeconômica, bem como os
documentos que deverão ser trazidos para comprová-la. Perguntas como renda
salarial, quantas pessoas contribuem com ela, números de bens móveis e imóveis,
situação habitacional, tipo de escola em que estudaram, se trabalham e há quanto
tempo, quanto gastam para vir à USP, o tempo de ida etc. Ainda há um espaço
para descreverem alguma particularidade não exposta nas perguntas que, obviamente,
receberá um parecer técnico.

Na classificação final da MORADIA, se vocês conseguiram uma pontuação grande, vocês
podem escolher entre a moradia no CRUSP ou um auxílio financeiro (bolsa moradia)
de R\$400,00 para vocês poderem alugar quartos, casas, hotéis, ou mesmo transporte
para ida e volta pra sua terra. Há casos em que a classificação final lhes
permite ter benefício apenas ao alojamento OU à bolsa, mas ao apartamento de
três quartos individuais (CRUSP) não.

Se vocês não conseguirem por nenhum desses meios, podem tentar a hospedagem, que é
simplesmente vocês ficarem no apartamento de alguém que more no CRUSP. Mas fiquem
atentos às datas de requerimento depois do resultado da seleção, pois vocês podem
ficar sem essa chance. Se mesmo assim vocês não conseguirem nada (bixes azarados!),
e acharem que os entenderam mal na entrevista ou coisa assim, vocês podem pedir para
entrar com recurso, e terem mais uma chance de esclarecer melhor a sua situação
(a.k.a. “cantar a assistente social”), ou também procurar a
AMORCRUSP (Associação dos Moradores do CRUSP que fica no Bloco F, das 14h às 18h).

{\bf Alimentação}

Vocês podem solicitar também o auxílio alimentação na página do SAS, que
consiste nos vale-bandex da USP e são válidos para almoço e jantar. Vocês deverão
passar por outra seleção que também inclui questionários sócio-econômicos,
comprovantes, e mais papéis.

{\bf Bolsa-trabalho}

Destina-se a alunos de graduação vinculados a projetos de extensão de serviços à
coletividade. Os projetos são selecionados anualmente, de acordo com sua relevância
para as finalidades da universidade pública e os estudantes vinculam-se por
afinidade acadêmica ou científica. Cada bolsa é de 1 (um) salário mínimo por
40 horas de trabalho mensais. Além da seleção socioeconômica feita pela
DPS (Divisão de Promoção Social), há uma seleção técnica feita pelos supervisores
dos projetos.

Mas fiquem espertos! Para tudo tem prazo e o SAS não é obrigado a ficar esperando a
boa vontade de aparecerem por lá de ninguém. Qualquer dúvida, liguem pro SAS.

{\bf Atendimento odontológico gratuito}

Antigamente para vocês agendarem o atendimento odontológico gratuito era necessário fazer uma
carteirinha no HU (Hospital Universitário) e em seguida, comparecer ao Bloco G
do CRUSP com a carteirinha e agendar, porém hoje em dia, graças à tecnologia,
isso não é mais necessário: você só vai precisar mandar uns e-mails, preencher uns
formulários e esperar \sout{uma eternidade}.

Primeiro você precisa mandar um e-mail solicitando triagem para triagemodonto.sas@usp.br,
vão te responder com um formulário que você vai preencher e enviar de volta, após isso eles
te colocarão na fila de espera pra triagem, no final deve demorar muitos meses pra
eles te chamarem e quando fizerem isso você só vai ter que comparecer ao Bloco G do CRUSP
munido do seu cartão USP (que a essa altura do campeonato você já estará com o definitivo).

No final será demorado, porém o serviço é gratuito e realmente ajuda. Eles também atendem casos
de emergência como por exemplo dor de dente, algum dente quebrado ou qualquer outra situação
do gênero. No caso de dúvidas, são poucos os que podem te ajudar nesse assunto, então o melhor
a se fazer é ligar pra lá: (11) 3091-3393.

%FIXME
\pagebreak
{\bf Setor de passe escolar}

Será nesse lugar que você resolverá boa parte dos seus problemas com o SAS e com o
Bilhete Único. Aqui é onde você carrega o vale bandex, entrega as documentações para
as bolsas, assina boa parte dos contratos e recebe orientação sobre as mais diversas
burocracias em que está se metendo.

Para as linhas da EMTU, SPTRANS, METRO, vocês devem entrar no site do SAS e fazer
o pré-cadastro para o respectivo cartão ({\tt http://sites.usp.br/sas/}). Talvez demore um pouco, pois a USP tem que
avisar para a SPTrans que vocês passaram na Fuvest. Fiquem atentos! Qualquer dúvida,
liguem para a seção de passe escolar do SAS: (11) 3091-3581.

Há também um cartão especial para alunos de universidades públicas provenientes
de outras cidades do interior de São Paulo. Se algum de vocês vem de algum desses domos
ignotos (como Resende, Caçapava ou Guaíra), pode se dirigir ao guichê da sua empresa de
transporte intermunicipal (Cometa, Danúbio Azul etc.), apresentar seu Cartão USP,
preencher um formulário e eles aguardarão confirmação do seu instituto. Então
você pega seu cartão na Seção de Alunos e, na compra de passagens entre São Paulo
e sua cidade-natal, paga 50\% do valor normal. Só não deixe isto por último na
sua lista de necessidades porque existe um período do ano em que os guichês
liberam seus formulários; em resumo, espiche suas orelhas e corra para a rodoviária.

\end{subsecao}

\begin{subsecao}{Atendimento Psicológico}

Existem inúmeros serviços na USP que oferecem atendimento psicológico
aos alunos, cada um com um foco e um público-alvo. Dentre eles, tem
o Serviço de Aconselhamento Psicológico (SAP), o Plantão Psicológico
LEFE e o Ateliê Aberto.

Bixes, sempre que precisarem, não tenham receio de procurar apoio
psicológico, por mais banal que vocês possam achar o motivo. Vocês
podem frequentar os plantões psicológicos para conversar sobre
qualquer coisa que esteja lhes incomodando, sejam professores,
matérias, EPs ou algo do tipo sem criar qualquer
problema para vocês no IME. Se for necessário, pode haver processos
de encaminhamento ou retorno para novas entrevistas.

Para maiores informações, acesse o site
\url{https://www5.usp.br/servicos/atendimento-psicologico/}

\end{subsecao}
%\pagebreak
% Utilize para começar uma nova página do lado esquerdo do Guia!

\begin{subsecao}{Bandejão}
%TODO Fazer as duas páginas ficarem na mesma "visualização" no guia impresso
%Ou seja, usar esquema Odd/Even do LaTeX.
% Isso foi deixado de lado nessa versão de 2016 para que o mapa dos circulares
% ficasse no lugar certo e não tivesse nenhuma página em branco
% (como o cleardoublepage ali em cima)
% 2017: não sei se entendi o que era para ser feito, mas usei o
%       \cleardoublepage para a imagem do calvin não ficar zuada.


Os bandejões, vulgarmente conhecidos como Restaurantes do SAS, são os lugares
em que vocês podem se alimentar razoavelmente a um preço analogamente razoável.
Os tickets custam R\$ 2,00, sendo estes carregados na carteirinha USP. O
lugar para carregar o ``bandejão único'' é no SAS, perto do bandejão
Central. Além disso, é possível verificar o saldo
atual e comprar mais créditos via boleto no site
{\tt https://uspdigital.usp.br/rucard/} ou pelo aplicativo "Cardápio USP",
disponível para Android e iOS (não se esqueça de olhar também os outros
aplicativos da USP disponíveis, pois podem ser muito úteis).
É uma ótima alternativa para evitar as infinitas filas dos guichês de venda. Os
créditos demoram até 02 (dois) dias úteis para estarem disponíveis.
Dica: Confira o saldo da carteirinha no site ou no aplicativo antes de tentar usar!

Existe uma lenda que sobremesas especiais aparecem em certos bandejões perto de
datas comemorativas. Mas não se engane - normalmente, elas não aparecem no
cardápio (talvez para não ficar muito cheio de gente).

\cleardoublepage
\figuragrande{bandex_calvin}

O cardápio semanal do bandejão pode ser visto no site {\tt
http://sites.usp.br/sas/} ou pelo aplicativo "Cardápio USP", mas se vocês
estiverem com preguiça de ver em um desses lugares, é
bem provável que um veterane já saiba e resolva informar se questionado com muita
educação.

O cardápio é geralmente composto de arroz (com opção normal e integral), feijão,
prato principal (carne/ovos), acompanhamento (legumes ou verduras refogadas,
cremes, molhos), salada (algumas folhas), sobremesa (frutas e/ou, às vezes, algo mais
requintado), pãozinho (normal e integral) e suco, além de temperos genéricos
(jamais perguntem do que eles são feitos). Se vocês são vegetarianos, o
acompanhamento nunca contém carne (e estão tentando fazer com que sejam veganos
também), e todos os bandejões têm uma alternativa vegetariana de PVT para o
prato principal, normalmente vindo na forma de ração, mas existindo as formas
de lasanha, kibe, e outros.


Não se esqueçam de levar as suas canecas do Kit-bixe se forem comer nos
bandejões, tanto para ostentar a posição de bixes do IME, quanto para salvar o
meio ambiente, evitando o uso de copos descartáveis. Além disso, levar a caneca
é extremamente vantajoso, já que a quantidade de suco que cabe no copo
descartável é pequena, fazendo vocês viajarem diversas vezes até a máquina de
suco.

PS: O Efeito Bandex é proporcional à quantidade de salitre utilizado em cada
bandejão.\\
PS 2: Nunca, em hipótese alguma, jamais, visite a cozinha do seu bandejão de
preferência, pois você corre o risco de nunca mais almoçar na vida. Como já
dizia o velho sábio ``A ignorância é uma virtude''.\\
PS 3: Playstation 3.

Consultem a tabela abaixo para decidir em qual dos bandejões vocês vão comer.

\figuragrande{bandex}
O horário das refeições é o seguinte:\\
Café da manhã: 7h às 8h30min (somente no Central, aos sábados fica aberto das
7h às 9h e aos domingos das 8h às 9h30)\\
Almoço: 11h15min às 14h15min (na Química é das 11h às 14h)\\
Jantar: 17h30 às 19h45 (não tem no PCO)\\

Para mais informações, visitem o site do SAS: {\tt http://sites.usp.br/sas/}

\end{subsecao}

\begin{subsecao}{Outros lugares para comer na USP}

Caso vocês, bixes, não queiram comer no bandejão, seja porque estão com medo do
prato do dia ou simplesmente porque querem comer algo diferente (e
possivelmente de sabor melhor), saibam que há alguns outros lugares na USP onde
vocês podem comer.

Alguns dos preços a seguir podem não estar 100\% corretos, mas a ordem de
grandeza sim =).

{\bf Lanchonete da Física}

Fica um pouco antes do bandejão da Física, nela vocês podem almoçar pagando R\$
40,90 o quilo, sempre tem uma razoável quantidade de saladas e pratos
quentes, além disso, tem churrasco com uma boa variedade de opções.
Além do almoço por quilo, lá vende alguns salgados e lanches como sanduíches e
beirutes. E um pouco antes da lanchonete, há um lugar que vende cookies por R\$
2,50 (3 unidades). Vocês vão sentir o cheiro dos cookies quando forem bandejar.

{\bf Restaurante do IPEN}

Fica na parte de cima da rua entre o IF e o Parque Esporte Para Todos, para
entrar lá vocês precisam da carteirinha USP. Lá o quilo é R\$ 24,50, porém não
tem uma variedade muito grande de comida, tem marmitex também que é R\$ 11,70.
Ele só abre para pessoas de fora às 13h e antes de ir para o restaurante vocês
devem se identificar.

{\bf Lanchonete da FAU}

Localizada dentro da FAU assim que você sobe a rampa. Lá além de lanches é
servido almoço tanto por quilo (custa R\$ 43,90) como prato feito. Tirando o por
quilo, qualquer outra coisa deve ser paga no caixa antes de retirar. Além
disso, no terceiro andar da FAU tem a mesinha de doces que os alunos deixam lá
para vender, a variedade e quantidade de doces varia bastante com o dia e o
preço deles costuma ser no máximo R\$2,50.

{\bf Restaurante da FEA}

Fica atrás da FEA, lá o quilo tem bastante coisa diferente e algumas até que
sofisticadas, porém o preço do quilo lá é R\$61,90, provavelmente é o mais caro
da USP. Geralmente é frequentado por professores e funcionários da USP. Além do
self-service, há algumas opções de prato feito que custam R\$22,00.

{\bf Trailers de lanche da \sout{ECA} Química}

Ficam na calçada com a Avenida Lineu Prestes, em frente ao Instituto de
Química. Nesses carrinhos vende pastel, sanduíches, churros, tapioca e salgados
em geral. Os preços dos lanches variam entre R\$3,00 e R\$10,00. As opções de
sanduíches e a quantidade de sabores de tapioca é bem grande e todos em geral
são muito bons. Já o pastel você pode escolher de 1 a 5 ingredientes para pôr
no pastel, portanto o preço varia de acordo com a quantidade de ingredientes que
vocês vão querer.

{\bf Cachorro-quente da Reitoria}

Fica na rotatória da biblioteca Brasiliana (ou a rotatória depois do CEPE, se assim preferir)
lá o cachorro-quente pode ser feito tanto no pão de cachorro-quente como na baguete. Além disso,
há a opção de pôr catupiry e de pôr uma salsicha adicional. O preço do cachorro-quente varia
entre R\$9,00 e R\$12,00.

{\bf Poke da Reitoria}

Recentemente abriu um lugar que vende Poke (um prato havaiano que mistura peixe cru em cubos, arroz, shoyu e frutas) ao lado do cachorro-quente. Esta iguaria custa aproximadamente R\$23, e você pode escolher os ingredientes que vão no seu prato, tipo Spoleto. Dá pra fazer um Poke vegano, com cogumelos em vez de peixe. Além de Poke, também vende Temakis, Hot Rolls, e até Açaí. A qualidade é boa mas as filas podem ser demoradas.

{\bf Pastel do IEE}

Esse lugar maravilhoso fica a apenas alguns passos do IME. Apenas às terças-feiras,
em um canto escondido do IEE (pergunte a um veterane como chegar!) fica o famoso pastel do IEE.
Há muitas opções de sabores, além de poder comprar um caldo de cana para acompanhar.
Geralmente, o preço do pastel é o mesmo ou mais barato que o pastel de feira.

\end{subsecao}

\end{secao}


% Um Pouco Sobre o DCE ---------------------------------------------------------
\input{sobre_o_dce.tex}

\quadrinhos{5}

% Hospital Universitário -------------------------------------------------------
\begin{secao}{Hospital Universitário}
   \begin{quote}\emph{O HU USP é o hospital de ensino de excelência utilizado
pelos Cursos de Atenção à Saúde da USP.  O hospital privilegia as pesquisas
relacionadas aos problemas de saúde  mais comuns da população brasileira. O
atendimento é regionalizado para o bairro do Butantã, sempre com enfoque no
ensino e pesquisa''}- Página do HU
   \end{quote}

O que isso quer dizer? Os estudantes de Medicina, Ciências Farmacêuticas,
Odontologia, Saúde Pública, da Escola de Enfermagem e do Instituto de
Psicologia, mantendo contato direto também com o Instituto de Ciências
Biomédicas, de Biologia, de Química, Faculdade de Arquitetura e Urbanismo,
Escola Politécnica, Escola de Comunicações e Artes, precisam de cobaias para
suas atividades/experiências. O HU é um santo lugar que recebe alguns fracos
de espírito que bebem demais e ficam incapacitados de fazer qualquer atividade
fisiológica. Vocês podem levar um comprovante de endereço e a
carteirinha USP para fazer o cartão do hospital. Assim vocês possuirão alguns
privilégios no atendimento, em situações de emergência vocês são atendidos
rapidamente (algo entre 600 minutos, como diz na senha de espera) e não
precisam soletrar o nome da mãe enquanto estiver desmaiado.

\end{secao}


% Tudo Que Vai Volta -----------------------------------------------------------
\begin{secao}{Tudo Que Vai Volta (até bixes)}

\begin{subsecao}{Ônibus}

Se vocês não têm como ir nem como voltar, temos algumas dicas:

\begin{enumerate}
  \item Trabalhem muito para comprar um carro,
  trabalhem mais para pagar a gasolina,
  venham para a USP de carro e, obrigatoriamente, deem carona a um veterane;

  \item Peçam a uma pessoa amiga para trazê-los e buscá-los durante seus
  longos anos de IME;

  \item Conheça alguém que, por sorte, mora perto da sua casa, estuda na USP,
  tenha o mesmo horário que você, seja legal e tenha carro. Traduzindo, s-o-n-h-e;

  \item Estiquem o dedão e esperem, esperem, esperem, esperem... a boa vontade
  de alguém para dar carona;

  \item Mudem-se para uma casa perto da USP;

  \item Peguem uma bicicleta e pedalem;

\end{enumerate}

OK, vocês decidiram ser bixes normais e pegar ônibus! Nesse caso, vocês
provavelmente estão em uma das seguintes opções:

\begin{enumerate}
  \item Vocês vêm de metrô pra USP;
  \item Vocês são super sortudos e acharam um ônibus que vai de casa pra USP.
\end{enumerate}

Se vocês vão usar ônibus para ir ou voltar da USP, vocês precisam conhecer
os quatro pontos de ônibus perto do IME: FAU I, Oceanografia, FAU II e FEA.
Tanto o ponto FAU I quanto o ponto da Oceanografia se localizam na Rua do Matão,
enquanto que os pontos FAU II e FEA ficam na Av. Prof. Luciano Gualberto (que a
partir de agora será chamada de Rua dos Bancos, para todo e todo o sempre).

Para tentar facilitar:
\begin{itemize}
	\item Av. Prof. Luciano Gualberto = Rua dos Bancos;
	\item Av. Prof. Lineu Prestes = Rua do HU;
	\item Av. Prof. Mello Moraes = Rua da Raia.
\end{itemize}

Cada um desses pares de pontos se localiza em lados opostos de suas respectivas
ruas. Se vocês chegaram na USP por um deles, o mais provável é que vão embora pelo
ponto do outro lado da rua.

Aqui está a lista de ônibus que passam em cada um dos pontos. Para mais detalhes
sobre cada linha, vocês podem usar o site da SPTrans ou o Google Maps.

\begin{subsubsecao}{Linhas municipais}

{\bf Ponto FAU I}

\begin{center}
	\begin{tabular}{|c|c|c|}
      \hline
	  Letreiro & Cor & Interligações (em ordem)\\
	  \hline
	  8022-10 - Metrô Butantã* & Laranja & Metrô: L4\\
      \hline
	\end{tabular}
\end{center}

{\bf Ponto da Oceanografia}

\begin{center}
	\begin{tabular}{|c|c|c|}
      \hline
	  Letreiro & Cor & Interligações (em ordem)\\
	  \hline
	  8012-10 - Metrô Butantã* & Laranja & Metrô: L4\\
      \hline
	\end{tabular}
\end{center}

{\bf Ponto FAU II}

(Esses os levam pra fora da USP)
\begin{center}
	\begin{tabular}{|c|c|c|}
      \hline
	  Letreiro & Cor & Interligações (em ordem)\\
	  \hline
	  177H-10 - Metrô Santana & Azul & Metrô: L4, L2, L3 e L1\\
	  7181-10 - Term. Princ. Isabel & Laranja & CPTM: L9\\
	  7411-10 - Praça da Sé & Laranja & Metrô: L4, L2, L3 e L1\\
	  7725-10 - Rio Pequeno & Laranja & - \\
	  809U-10 - Metrô Barra Funda & Laranja & L2 \\
	  8022-10 - Metrô Butantã* & Laranja & Metrô: L4\\
      \hline
	\end{tabular}
\end{center}

Esses dois ônibus passam no ponto, mas estão CHEGANDO na USP, portanto não sejam
ridículos dando sinal para eles pararem e muito menos peguem um desses ônibus
nesse ponto.

\begin{center}
	\begin{tabular}{|c|c|}
	  \hline
	  Letreiro & Cor\\
	  \hline
	  701U-10 - Butantã-USP & Azul\\
	  702U-10 - Butantã-USP & Laranja\\
	  \hline
	\end{tabular}
\end{center}

{\bf Ponto FEA}

(Esses os levam pra fora da USP)
\begin{center}
	\begin{tabular}{|c|c|c|}
      \hline
	  Letreiro & Cor & Interligações (em ordem)\\
	  \hline
	  701U-10 - Metrô Santana & Azul & Metrô: L4, L2, L3 e L1\\
	  702U-10 - Term. Pq. D.Pedro II & Laranja & Metrô: L4, L2 e L3\\
	  7725-10 - Terminal Lapa & Laranja & CPTM: L8\\
	  8012-10 - Metrô Butantã* & Laranja & Metrô: L4\\
      \hline
	\end{tabular}
\end{center}

Novamente, esses quatro ônibus passam no ponto, mas estão CHEGANDO na USP!
\begin{center}
	\begin{tabular}{|c|c|}
	  \hline
	  Letreiro & Cor\\
	  \hline
	  177H-10 - Butantã-USP & Azul\\
	  7181-10 - Cidade Universitária & Laranja\\
	  7411-10 - Cidade Universitária & Laranja\\
	  809U-10 - Cidade Universitária & Laranja\\
	  \hline
	\end{tabular}
\end{center}

As linhas marcadas com um * são as linhas circulares da SPTrans. Segue abaixo suas descrições!

\end{subsubsecao}

\begin{subsubsecao}{Circulares}

Também conhecido como ``circulenda'' ou ``secular'' (aos sábados, ``milenar'',
e aos domingos, ``anos-luz''), é o meio de transporte mais barato dentro da USP.
Foi criado para os USPianos se locomoverem dentro do Campus, mas em muitas vezes
é melhor andar do que ficar esperando. Existem 2 itinerários distintos, com
trajetos aproximadamente reversos. Fiquem atentos para não darem uma de bixe
burro e se perderem, hein?

Em 2012 foram implantadas as linhas 8012/10 (circular 1) e 8022/10 (circular 2)
- Metrô Butantã/Cidade Universitária,
que funcionam como circulares USP, e nós alunos não pagamos, pois elas aceitam
o bilhete USP (BUSP - Retire logo o seu!!).

Para acompanharem as rotas dos circulares e acompanhá-los em tempo real, bem como
ver os portões do campus que estão abertos e seu horário de funcionamento, pode-se
usar o aplicativo “Portões USP", disponível para Android e escrito por ex-alunos
do IME.

Para os bixes que não gostam de ler (preferem imagens) colocamos o
mapa das duas linhas nas próximas páginas.

\mapa{8012-10}
\mapa{8022-10}

\end{subsubsecao}

\begin{subsubsecao}{Como ir e vir do IME pelo metrô Butantã}

Como sabemos que vocês, bixes, chegam muito perdidos, e muitos pularam todas essas informações
sobre os pontos de ônibus e por isso podem não saber como fazer seu trajeto, resolvemos colocar
aqui um dos trajetos mais comuns que boa parte de vocês vão usar.

Para chegar no IME a partir do metrô Butantã, vocês devem pegar o circular 1 (8012-10).
Se vocês pegarem o circular 2 (8022-10), em algum momento vocês vão chegar no IME, mas
o circular 1 é muito mais rápido.

O circular 1 vai fazer o seguinte trajeto: Faculdade de Educação, CEPEUSP, Praça do Relógio
Solar, Biblioteca Brasiliana e Faculdade de Economia, Administração e Contabilidade. A partir
daí, deverão descer no ponto da FEA.

Quando estiverem indo embora, peguem o circular 2 (8022-10) no ponto da Oceanografia. Se vocês
pegarem o circular 1 (8012-10) no ponto FAU I, também vão (em algum momento) chegar no metrô
Butantã, mas boatos dizem que o circular 2 vai mais rápido. Aí, basta ficar no ônibus sentado
(se conseguir um lugar!) até chegar no metrô Butantã, onde todos do ônibus vão descer.

\end{subsubsecao}

\begin{subsubsecao}{Linhas intermunicipais}

Agora, se vocês moram mais longe ainda (outra cidade, outro estado, outro país,
outra dimensão...) e não querem ou não podem se mudar para São Paulo, existem
algumas linhas de ônibus fretados para cidades mais próximas (ou não). Se por
acaso alguma cidade não esteja aí, procurem se informar a respeito, pois não significa
necessariamente que não haja ônibus da USP para lá. Aí estão elas:

\begin{itemize}
  \item {\bf Empresa Urubupungá.}\\
    Tel: 0800 11 4777
    Site: {\tt www.urubupunga.com.br}\\
    280BI1- São Bernardo do Campo (Centro)\\
    Cor: Cinza\\
    Onde pegar para sair da USP: ponto FAU II\\
    Av Magalhães de Castro, Av Marginal Pinheiros (Shopping Eldorado), Av Dos
    Bandeirantes, Av Eng. Luiz Carlos Berrini, Av Roque Petroni Jr. (Shopping
    Morumbi), Av Prof Vicente Rão, Av Cupecê (Diadema), Av Fábio Eduardo Ramos
    Esquivel (Diadema), Av Piraporinha (Diadema), Av Lucas Nogueira Garcez
    (Diadema), Av Urubupungá.

  \item {\bf Fretados Jundiaí - USP}\\
    Viação MIMO\\
    Tel: (11) 4606-8222\\
    {\tt www.viacaomimo.com.br}\\
    Principais horários:\\
    Ida: 6h00; 6h30; 7h00; 12h50; 18h00 (na Rodoviária de Jundiaí)\\
    Volta: 11h40; 15h00; 16h40; 17h05; 18h10; 21h00; 22h49 (No ponto FEA)

  \item {\bf Bragança}\\
    N.S. Fátima\\
    {\tt www.saexbra.com.br}\\
    Tel: 4032-4723 e 7344-2007\\
    Ida: 05h50 (Lgo do Tabão/Habbibs)\\
    Volta:17h00 (Av Prof. Almeida Prado)

  \item {\bf Campinas}\\
    Sta. Cruz\\
    Tel: 3868-5995\\
    {\tt www.gruposantacruz.com.br}

  \item {\bf Santos}\\
    Náutica Turismo\\
    Tel: (13) 9112-8860

  \item {\bf Santos / S. Vicente}\\
    Transul\\
    Tel: 6954-4466

  \item {\bf ABC}\\
    Dinâmica ABC-USP\\
    4352-0565 / 4109-0172

  \item {\bf Sorocaba}\\
    Fretado Diurno\\
    (15) 9715-1676 (Márcio)

\end{itemize}

\end{subsubsecao}

\end{subsecao}


\begin{subsecao}{Veículos no Campus}

Saibam por onde entrar na USP. Lembrem-se de terem sempre a carteirinha USP ou
comprovante de matrícula com RG em mãos, pois pode ser solicitado principalmente
nos horários de entrada controlada.
\begin{itemize}
  \item {\bf Portaria 1 (P1):} R. Afrânio Peixoto. Funciona 24h por dia todos os
    dias, mas a entrada é controlada de segunda à sexta das 20h às 05h, aos sábados
    após as 14h, aos domingos e feriados o dia inteiro. É por onde entram os
    ônibus municipais.

  \item {\bf Portaria 2 (P2):} Av. Escola Politécnica. De segunda a sexta, acesso
  livre das 5h às 21h e controlado das 21h às 24h. Fechada aos sábados, domingos
  e feriados. Única entrada para caminhões.

  \item {\bf Portaria 3 (P3):} Av. Corifeu de Azevedo Marques. Tem o mesmo horário
    de funcionamento do P2, porém abre aos sábados até das 5h às 14h e tem acesso
    controlado das 14h às 25h.

  \item {\bf Portaria P' (Plinha):} R. Eng. Teixeira Soares. Funciona de segunda a
  sexta das 6h às 21h. Fechado de sábados, domingos e feriados.


  \item {\bf Portarias de pedestre (Mercadinho, São Remo, HU, CPTM e
      Vila Indiana):} Funcionam de segunda a sábado, das 5h às 20 h. Acesso controlado das 20h
      às 24h. Nas portarias do Mercadinho e Vila Indiana, pode entrar entre 0h e 5h, além
      de poder entrar em domingos e feriados, mas o acesso é controlado.

\end{itemize}

Os ônibus entram no sábado até as 14h e não entram no domingo, com exceção aos circulares, que entram na Universidade a qualquer hora. Se vocês vêm de carro, saibam que a universidade dispõe de bolsões de estacionamento
gratuito em torno das Unidades.

\end{subsecao}

\begin{subsecao}{Pontos de táxi}
Existem alguns pontos de táxi espalhados pela Cidade Universitária. As frotas
operam de segunda a sexta-feira,das 7h às 23h, além de aos sábados, das 7h às 17h.
Eis suas localidades:

\begin{itemize}
\item Praça das Agências Bancárias\\
Fone: (11) 3091-4488

\item Praça da Reitoria\\
Fone: (11) 3091-3556

\item Hospital Universitário\\
Fone: (11) 3091-3536
\end{itemize}
\end{subsecao}

\end{secao}


% Guia de jogos da Vivência ----------------------------------------------------
\input{guia_de_jogos.tex}

%FIXME GAMBIARRA para não quebrar página num lugar zuado
\pagebreak

% Dicas ------------------------------------------------------------------------
\begin{secao}{Dicas}

Como todos os bixes chegam perdidos, aqui vão algumas dicas pra vocês não ficarem perguntando o tempo todo:

%TODO Checar isso aqui.
{\bf Bancos e Caixas Eletrônicos:} agências do Santander, Bradesco,
Banco do Brasil, Caixa e Itaú na Av. Prof. Luciano
Gualberto; caixas eletrônicos perto do bandejão da Física, em frente à reitoria,
no CEPE, enfim, em lugares aleatórios.

{\bf MatrUSP:} é um site criado por alunos do BCC para simular grades horárias
para matrículas de disciplinas. Assim, vocês conseguem saber direitinho se suas
disciplinas vão coincidir e quanto tempo livre vocês vão ter para ficar na
vivência entre as aulas. Acessem: \url{http://bcc.ime.usp.br/matrusp}

{\bf USPAvalia:} outro site criado por alunos do BCC (pois é, olha só!) que
contém inúmeras avaliações da comunidade sobre disciplinas oferecidas e
seus respectivos professores. Ótimo para saber se vale a pena ou não pegar
essa ou aquela turma de uma dada disciplina, e vocês mesmos podem contribuir com
suas próprias avaliações e comentários! Acessem: \url{https://uspavalia.com}

\begin{subsecao}{Cultura na USP}

Vale aqui lembrar que a entrada em vários desses museus é gratuita para alunos da
USP, então vale muito a pena visitar!

{\bf Museu de Arqueologia e Etnologia (MAE):} ao lado da Prefeitura do Campus,
possui um dos maiores acervos de artefatos arqueológicos e etnográficos do Brasil.

{\bf Museu de Arte Contemporânea (MAC):} próximo ao CRUSP, nele são expostas produções
artísticas nacionais e estrangeiras.

{\bf Museu do Brinquedo:} fica na Faculdade de Educação, Bloco B. Seu acervo conta 
com itens datados do início do século XX, incluindo brinquedos, jogos, materiais
pedagógicos e um acervo fotográfico. Tem a exposição "Cenas Infantis", que fica na
bilbioteca da Faculdade de Educação.

{\bf Museu do Crime da Polícia Civil:} na Academia de Polícia perto do P1. Seu acervo
reúne ferramentas, objetos e documentos utilizados em delitos de grande repercussão, 
além da história de criminosos cujos atos ficaram marcados na imprensa e na sociedade
brasileira.

{\bf Museu do Instituto Oceanográfico:} adivinha? Esse museu mantém uma exposição voltada
a dinâmica, a estrutura e a biodiversidade dos oceanos.

{\bf Museu da Geociências:} Lá mesmo. Conta com amostras de rochas, gemas, meteoritos 
e fósseis.

{\bf Instituto Butantan:} Próximo à História. É um dos maiores acervos de pesquisa
biológica do mundo, conta com a presença de diversas cobrinhas.

{\bf Museu de Anatomia Veterinária:} Perto do P3 (portão da Corifeu). Seu acervo possui
uma coleção de dados e fotos de esqueletos, além de modelos anatômicos e animais
preservados.

{\bf Museu Paulista, vulgo Ipiranga:} fora da USP, no Parque da
Independência - S/N  - Ipiranga. Atualmente está fechado para visitação,
pois está sendo restaurado.

{\bf Museu de Zoologia:} também fora da USP, na Av. Nazaré, 481  -
Ipiranga. Sua exposição abriga uma série de animais empalhados, fósseis,
réplica de fósseis etc.

{\bf CinUSP Paulo Emílio:} Dentro do campus, próximo ao bandejão central, existe uma sala de cinema. Durante todo o ano ocorrem várias mostras cinematográficas, nas quais são exibidos inúmeros filmes. As sessões são gratuitas e a programação pode ser conferida no seguinte site: {\tt http://www.usp.br/cinusp/}


\end{subsecao}

\begin{subsecao}{Onde beber?}

Se vocês são dos bixes que curtem entornar os canecos de vez em quando, então
agora devem estar pensando, até que enfim vamos falar de algo que presta!
Lembre-se de levar um veterane para pagar a ele algumas doses, pois é graças a
eles que vocês estão recebendo essas dicas, então sem mais delongas, aqui vão
alguns lugares firmeza para se fazer isso à vontade:

{\bf FAU:} Na vivência da FAU sempre vende cervejas. Para chegar, basta ir no 
primeiro andar à direita até o fim.

{\bf Física:} Embora seja a Física, lá é um lugar gostoso para comer alguns
salgados na lanchonete e tomar algumas cervejas na atlética.

{\bf FFLCH:} Vá até o prédio da História e Geografia e vá até o Aquário, fica 
logo à esquerda na entrada do vão, se estiver alguém lá dentro eles vendem cerveja.

{\bf FEA:} Entrando na FEA, ande até o final, siga reto depois da biblioteca,
passe por um portãozinho, vire à esquerda e tem uma entrada antes do restaurante.
Lá é a vivência da FEA onde vendem as cervejas geladinhas.

%FIXME Qual é a situação atual do QiB?
{\bf ECA:} Famosa Quinta i Breja, adivinha que dia da semana isso acontece?
Acertou, toda quinta-feira na prainha da ECA a partir das 19h!

{\bf Rei das batidas:} Muito famoso não só por quem estuda na USP, o Rei,
como é carinhosamente chamado, fica fora da USP, saindo pela P1. Vende
diversas batidas e, é claro, cerveja...

{\bf Bar do frango:} Não se sabe qual é o verdadeiro nome desse bar, mas ele é
uma alternativa ao Rei, quando este se encontra muito lotado. Apesar do péssimo atendimento e do aspecto horrível do lugar é bom para beber sem
tumulto. É frequentado principalmente no começo do ano. Se encontra atrás do
Rei.

É claro que também há as festas, que ocorrem em qualquer lugar da USP e quase todas as sextas tem, e nelas
há ainda outras misturas alcoólicas impossíveis.

É nosso dever informar também que o álcool é uma substância altamente viciante
e, quando bebida em excesso, pode trazer graves consequências à sua saúde, às
vezes à saúde de outra pessoa, à sua família e principalmente ao seu bolso.
Lembrem-se também que é proibida a venda de bebidas alcoólicas dentro da USP, então
não saiam da vivência dos respectivos lugares com a latinha na mão.

\end{subsecao}
\end{secao}

\quadrinhos{9}

%FIXME GAMBIARRA para não quebrar página num lugar zuado
\pagebreak

% Melodias para a bixarada -----------------------------------------------------
%\pagebreak
\begin{secao}{Músicas para a Bixarada}

\begin{subsecao}{Hino Cabeção}
\begin{verse}

Ó IME USP, meu amor\\
Para sempre campeão\\
Cálculo só me traz dor\\
E a REC emoção

Com a cachaça na mão\\
Para teus times vou torcer\\
Jogar com raça e coração\\
Vermelho e branco até morrer

ôôô ôôÔÔ ôôÔÔ ôôÔÔ

Eu vim aqui\\
Só pra gritar\\
Dá-lhe! IME! USP!
\end{verse}
\end{subsecao}

\begin{subsecao}{Musiquinhas de inter}
\begin{verse}

"O IME é o IME do IME"

"IME chupa IME lambe"

"IME ao contrário é M. M de Mackenzie"

"O IME é uma droga, vamos fumar o IME"

\end{verse}
\end{subsecao}

\begin{subsecao}{Caboclo da MAT}

{\em Cantar como ``Faroeste Caboclo'' da Legião Urbana. Essa música é em memória 
à FUVEST, quando Matemática Aplicada e POLI pertenciam à mesma carreira}

\begin{verse}
\footnotesize{
Cheio de medo em setembro Joãozinho viu que seus dedos tremiam pra fazer a
inscrição\\
Deixou pra trás a namorada, a motoca, o futebol e as festinhas pra rachar na
revisão\\
Quando criança só pensava em ser engenheiro ainda mais com o dinheiro que
sonhava em ter na mão\\
Era o CD lá do colégio onde estudava e todo mundo admirava o boletim desse
cuzão\\
Ia pra igreja só pra rezar pro seu santo pra pedir a sua ajuda pra prestar
vestibular\\
Sabia mesmo que ia ser barra pesada porque tinha muito japa pra tomar o seu
lugar\\
O ano todo se propôs a estudar, passava o dia sem ligar a televisão\\
Nos feriados não ia viajar, ficava em casa treinando redação\\
Fazia todos os exercícios da apostila e no fim de cada aula ia falar com o
professor\\
Às quinze horas ia pro laboratório ver as mitocôndrias da aula anterior\\
Não entendia como o militarismo dominou nosso país por vinte anos de terror\\
Ficou cansado de tentar achar resposta e desceu pra lanchonete pra afogar a sua
dor\\
E lá chegando foi tomar um cafezinho e encontrou um concorrente com quem foi
falar\\
E o concorrente aumentou seu desespero pois manjava muita coisa que ele tinha
que estudar\\
Dizia ele, eu vou prestar o ITA... Nesse país prova pior não há\\
E se não der eu vou pegar engenharia, lá na POLI eu vou tomar o seu lugar\\
E João não gostou dessa proposta, ele disse ``ai que bosta, eu tô passando mal''\\
Ele ficou bestificado com a idéia de pegar lista de espera só depois do carnaval\\
Meu Deus, é pior ainda, no ano novo eu posso estar lá na Mauá\\
É brincadeira querer ser engenheiro e só descolar emprego em Taguatinga\\
Na sexta-feira ele morria de vontade de correr pro banheiro se borrando de pavor\\
E conhecia muita gente arrogante que passava do seu lado se dizendo um terror\\
Ele estudava o relevo da Bolívia, função quadrática e modular\\
E nos domingos então ele fazia tarefa mínima e complementar\\
E Joãozinho até a morte se esforçava e o tempo mal sobrava pr'ele se alimentar\\
E via às duas horas o Vestibulando que passava todas as dicas sobre o vestibular\\
Mas ele não queria mais conversa e decidiu que em novembro era hora de rachar\\
Ele pirou que precisava estudar tanto, virou um bitolado e começou a delirar\\
E logo, logo os malucos da sua idade viram a calamidade, tem babaca novo aí\\
E o nosso Joãozinho ficou louco e bateu em todos os japoneses dali\\
Seus amigos preocupados com a sua sorte deram uma fita de rock pr'ele relaxar\\
Mas de repente sob uma má influência dos boyzinhos lá do fundo começou a zoar\\
Já na primeira fase ele penou e só passou porque o corte foi sessenta e três\\
A demência tomou a sua mente : ``Vocês vão ver, eu vou pegar vocês!!!''\\
Agora Joãzinho era fodido e estava decidido que não ia se dar mal\\
Sacava toda a trigonometria e manjava de limites, derivada e integral\\
Foi quando conheceu uma menina e de toda aquela zona ele se arrependeu\\
Maria Lúcia era uma bitola linda e o coração dele pra ela o Joãzinho prometeu\\
Ele dizia que devia estudar, pois engenheiro ele queria ser\\
Maria Lúcia, pra sempre vou te amar, Engenharia com você quero fazer\\
O tempo passa e um dia chega a hora de fazer segunda fase coitadinho do João\\
E ele faz uma prova perigosa diz que espera uma resposta, pode ser um sim ou não\\
Não vou correndo pra banca de jornal nem pra pátio do cursinho isso eu não faço
não\\
Pois eu prefiro ficar na minha casa esperando o resultado com o cu na mão\\
Maria Lúcia vai comprar o tal jornal e logo após achar seu nome ela procura o de
João\\
Mas ela volta com tristeza no olhar, olha pra ele e diz ``você pegou a quarta
opção''\\
Você passou na sua quarta opção, você passou na sua quarta opção\\
Bacharelado em Matemática é um tesão, eu vou sofrer as conseqüências como um cão\\
Não é que Joãozinho estava certo, seu futuro era incerto mas foi se matricular\\
Matriculou-se e no meio da zoeira descobriu que tinha muitos como ele no lugar\\
Fez inscrição pro remanejamento e talvez no fim do ano transferência ia tentar\\
E Joãozinho mantinha a esperança de um dia ir pra Poli estudar química\\
Mas acontece que um tal Professorzzini terrorista de renome apareceu por lá\\
Ficou sabendo dos planos de Joãozinho e decidiu que com suas notas ele ia se
ferrar\\
E ele teve que largar cálculo dois mesmo sabendo derivar e integrar\\
E decidiu deixar estat pra depois que o Moretin voltasse a lecionar\\
Professorzzini, professor mais sem vergonha com sua prova enfadonha fez todo
mundo dançar\\
Desvirginava bixetes inocentes e o nabo era tão quente que nem dava pra sentar\\
E Joãozinho há muito não via sua amada, e a saudade começou a apertar\\
Eu vou pra Poli eu vou ver Maria Lúcia, já está em tempo de a gente se encontrar\\
Chegando à Poli então ele chorou quando viu Maria Lúcia namorando um japonês\\
Oh, Maria Lúcia, quanto que você mudou, que estrago que a Poli te fez\\
Joãozinho era só ódio por dentro e então o japonês para um duelo ele chamou\\
Amanhã às duas horas no biênio, ou na praça do relógio, seja lá onde for\\
E você pode escolher as suas armas: derivadas ou matrizes de qualquer versor\\
Que eu provo que o sub-espaço nulo é o coração dessa piranha a quem jurei o meu
amor\\
E Joãozinho não sabia o que fazer quando escutou um papo lá no bandejão\\
Onde falavam dum duelo que iam ver dizendo a hora, o local e a razão\\
No sábado então às duas horas toda a Poli sem demora foi lá só pra assistir\\
Um japa que botava pelas costas, encoxou Maria Lúcia e começou a sorrir\\
Sentindo um ódio na garganta João olhou pros cabacinhos e pros trouxas a
aplaudir\\
E olhou pros pipoqueiros e as bancas de cachorro-quente que passavam por ali\\
E se lembrou de quando era uma criança e de tudo que vivera até ali\\
E decidiu entrar de vez naquela dança, se a Poli é um circo, e daí\\
E nisso o céu abriu seus olhos e então Maria Lúcia ele reconheceu\\
Ela queria fazer Álgebra dois pra provar que a Poli não a emburreceu\\
Politécnico, eu sou homem coisa que você não é, e não me contento em por nas
costas não\\
Some daqui filha da puta sem vergonha vai pra casa tocar bronha o seu destino é
ser bundão\\
E Joãozinho deu as costas para os dois, foi pra pura onde encontrou o seu valor\\
Maria Lúcia se arrependeu depois prestou Fuvest mas no IME não entrou\\
E a todos declarava que o nosso Joãozinho era gênio que escapou de se foder\\
Que na alta burguesia lá da Poli todo mundo é bunda mole ninguém sabe o que
fazer\\
E foi dar monitoria no cursinho pra avisar aos molequinhos pra não esquecer\\
Ele queria era avisar toda essa gente engenharia é pra demente que só quer
sofrer.\\
}
\end{verse}
\end{subsecao}

\end{secao}


% Utilidades -------------------------------------------------------------------
\begin{secao}{Utilidades}

\begin{subsecao}{Na WEB}

{\tt www.usp.br} - Página da USP. Aqui vocês encontrarão notícias e eventos da
universidade, bem como informações gerais.

{\tt www.ime.usp.br} - Página do IME.
Nesse link, vocês poderão ver detalhes sobre os cursos e obter informações sobre
a faculdade.

{\tt uspdigital.usp.br/jupiterweb} - Sistema JúpiterWeb. Aqui vocês vão
encontrar a grade horária e, mais tarde, vocês poderão fazer as matrículas nas
matérias que irão cursar no semestre, além de ter acesso às suas notas.

{\tt paca.ime.usp.br} - É bixes... acham que vai ser essa moleza pra sempre? Se
vocês acham, estão muito enganados! Daqui a pouco vocês vão receber uma senha para
poder enviar seus EPs (vide glossário) nesse endereço... (e não adianta fazerem
chantagem emocional que o Paca só vai aceitar até 23h55... não entendeu? vocês
vão entender...)

{\tt fb.com/Segredos-IME-USP-403785843127530} - Precisa desabafar?
Precisa de conselhos ou opiniões? O Segredos IME - USP é o lugar perfeito para
tudo isso e muito +.


{\tt fb.com/SpottedImeUsp} - Viu alguém interessante? Quer mandar uma cantada
pro crush mas tem vergonha? Manda um spotted! Afinal, só o x deve ficar isolado.

% {\tt www.xkcd.com} - Webcomic sobre matemática e computação. Origem de muitas
% piadas que vocês ouvirão por aí.

E por fim..

{\tt www.google.com.br} - Tudo o que vocês precisam tá no google. Se não estiver,
então não existe. Aproveitem e procurem o Código de Ética da USP (ou baixem ele em
\url{http://www.prg.usp.br/wp-content/uploads/CodigoEtica.pdf})

\end{subsecao}

\begin{subsecao}{Telefones}
Colocamos essa seção aqui apenas para que vocês, bixes, não nos encham com perguntas
sobre passes e plantão:

{\bf SAS:} Seção de Passe Escolar: {\tt 3091-3581}

{\bf Seção de Alunos:} {\tt 3091-6149}

{\bf Telefonista da USP:} {\tt 3091-4313}

{\bf Disk Trote:} {\tt 0800-012-090} -- \underline{Não deixe de ligar se você sofrer algum trote violento!!}

{\bf Plantão de Cálculo e Álgebra Linear (serviço gratuito):} {\tt 3037-1773}
%Número do orelhão do IME - ou seja, existe uma grande probabilidade de ser atendido ^^

\end{subsecao}
\end{secao}


% Glossário --------------------------------------------------------------------
\begin{secao}{Glossário}

Esta parte é, sem dúvida, uma das mais importantes do guia. Aqui vocês
encontrarão todas as explicações para as maiores dúvidas do universo, apesar de
já sabermos que a resposta para a vida, o universo e tudo mais é 42. Com
certeza, após lerem este trecho sua vida vai mudar: vocês saberão, por exemplo,
porque o céu é azul e com quantos paus se faz uma canoa.

\begin{subsecao}{Sobre as matérias}

{\bf Teorema:} Um teorema é uma afirmação que pode ser provada. Provar
teoremas é a principal atividade dos matemáticos. Deles surgem Lemas,
Corolários, Proposições, e tantas outras coisas que você só vai entender
completamente o que significam quando precisar escrever sobre eles, o que vai
acontecer logo, logo!

{\bf Iniciação Científica:} grupo de estudantes, coordenado por um professor, que
estuda um determinado assunto paralelamente ao curso. No caso de estudantes que
queiram bolsas de estudo, é adotado um plano de estudos mais rigoroso.

{\bf P1, P2, P3:} sofrimento de todo aluno de graduação. A quantidade de sofrimento
depende do professor da disciplina.

{\bf SUB:} prova que você faz quando vai mal nas P1, P2 e P3, ou
quando você simplesmente não vai. Sua aplicação e utilidade depende do
professor ministrante.

{\bf REC:} prova que você faz quando vai mal na SUB. Pra fazer a REC, precisa ter
pelo menos média 3 nas outras avaliações.

{\bf DP:} matéria que você faz quando vai mal na REC.
\end{subsecao}

\begin{subsecao}{Sobre programas}

{\bf Computador:} objeto com vontade própria, sensível, que requer muito
carinho e atenção.

{\bf EP:} Exercício-Programa. Algo que você vai ter que fazer muitas vezes, e
vai dar muito trabalho.

{\bf GCC:} compilador mais recomendado para seus EP's, por suas inúmeras
qualidades. Atenção: ele ainda fará você se sentir incompetente.

{\bf Hello World:} Um clássico da programação universal.

{\bf Segmentation Fault:} Efeito computacional aleatório causado pela ``véspera
de entrega de EP''. Desenvolvido por Murphy.

{\bf Stack Overflow:} mensagem que aparece na tela do computador
quando ele se recusa a funcionar.

{\bf Teorema Fundamental do EP:} ''O EP só funcionará no dia da entrega.'' Não
confunda com o Corolário 42 da Lei de Murphy: ''O EP só {\bf não} funcionará no
dia da entrega!''

{\bf Linux:} Sistema operacional criado totalmente em linguagem C, graças a um
esforço mundial de milhares de programadores e experts em informática, composto
por aproximadamente 7 mil arquivos e 5 milhões de linhas, e com o qual você não
tem capacidade para trabalhar.

{\bf Windows:} Vírus. Porém tão bem mascarado que parece até a coisa correta a
se usar.
\end{subsecao}

\begin{subsecao}{Sobre a USP}

{\bf CEPE:} Centro de Práticas Esportivas - lugar onde você poderá praticar
todos os esportes que quiser.

{\bf SAS:} ao lado da praça do relógio. É onde os alunos fazem a carteirinha
de passes de ônibus, EMTU e metrô, além de solicitar os auxílios eteceteras que
estão explicados na seção respectiva.

{\bf CRUSP:} Conjunto residencial da USP. Se você se inscrever, torça para
pegar um apartamento num dos blocos já reformados, ou então torça para não
conseguir nenhum.

% {\bf Colmeia:} conjunto de favos.

% {\bf Favos:} um monte de prédios hexagonais encravados no meio do CRUSP.

{\bf CINUSP:} Cinema da USP localizado no favo 4 da Colmeia. Toda semana ele
passa um filme de qualidade. Informe-se sobre a programação em {\tt www.usp.br/cinusp}.

{\bf Pelletron:} prédio da Física que na realidade é um acelerador de
partículas.

{\bf P1, P2, P3:} não confundam com as do item das matérias. Aquelas eras AS P1, P2, P3;
estes aqui são OS P1, P2, P3. São os portões da USP. O P1, que é o principal, se localiza
no cruzamento da Av. Afrânio Peixoto com a Rua Alvarenga, além de que é por ele que
entram os circulares vindos do metrô. O P2 fica na Av. Escola Politécnica. Já o P3
fica na Av. Corifeu de Azevedo Marques. Fique atento quanto aos horários dos portões,
pois eles não ficam abertos o tempo todo.

{\bf H.U. (Hospital Universitário):} para onde são mandados bixes que se machucam na Semana de Recepção.
Lá são realizados os tratamentos e experiências com
exposição a radiação, exposição a aspirantes a médicos, teste de
paciência/resistência a dor assim que pega a senha, alto nível de gesso no
estômago, queda espontânea (ou não) de cabelo etc.

{\bf Vet (Veterinária):} para onde são mandados biches que se machucaram na Semana de Recepção.

{\bf Psico:} Para onde são mandados os bixes (e veteranes também) que querem plantões psicológicos porque saúde mental é importante deve ser tratada com seriedade, Porra.

\end{subsecao}
\end{secao}



% Considerações Finais ---------------------------------------------------------
\begin{secao}{Considerações Finais}

Este guia chegou ao fim. Vocês devem estar pensando ``E agora? O que nós, bixes,
vamos fazer, sozinhos nessa faculdade?''.

Se precisarem de alguma ajuda, basta procurarem algum veterane que ele te
ajudará (provavelmente ele também não saberá a resposta, mas talvez conheça alguém que saiba).

Não esqueçam de comprar seu kit-bixe para ajudar a recepção e tornar a recepção
do ano que vem mais legal que a de vocês. Participem da Semana de Recepção,
carinhosamente preparada pelo Instituto e pelos seus veteranes devidamente
identificados. Último conselho: guardem este guia para o resto de sua graduação.
Por mais ``bixes'' que sejam, vocês precisarão dele.

Qualquer sugestão, elogio, presente, bajulação e outras coisas LEGAIS para a
Comissão, pode ser enviada em {\tt fb.com/recepcaoimeusp}

Se vocês encontraram algum erro dos vários espalhados como exercício pelo texto,
nos enviem uma mensagem no Facebook ou clonem nosso repositório em
\url{https://github.com/kazuo256/guiadobixo} e ajudem vocês também na próxima
edição do Guia do bixe!

\end{secao}


%FIXME
%REFTIME
% O guia precisa ter um número de páginas totais múltiplo de 4, e a última
% página tem que estar vazia

\newpage
\thispagestyle{empty}
\pagebreak

\newpage\phantom{skip}
\thispagestyle{empty}
\pagebreak

\newpage\phantom{skip}
\thispagestyle{empty}
\pagebreak

\newpage\phantom{skip}
\thispagestyle{empty}
\pagebreak

\end{document}
